%!TEX root = forallxyyc.tex
\part{Dedução natural para a LVF}
\label{ch.NDTFL}
\addtocontents{toc}{\protect\mbox{}\protect\hrulefill\par}

\chapter{A própria ideia da dedução natural}\label{s:NDVeryIdea}

Bem antes, em \cref{s:Valid}, dissemos que um argumento é válido \ifeff{} não há caso em que todas as premissas sejam verdadeiras e a conclusão seja falsa.

No caso da LVF, isso nos levou a desenvolver tabelas-verdade. Cada linha de uma tabela-verdade completa corresponde a uma valoração. Assim, diante de um argumento da LVF, temos uma maneira muito direta de avaliar se existe uma valoração na qual as premissas sejam verdadeiras e a conclusão seja falsa: basta percorrer exaustivamente a tabela-verdade.

No entanto, as tabelas-verdade talvez não nos deem muita \emph{compreensão}. Considere dois argumentos da LVF:
\begin{align*}
P \eor Q, \enot P & \therefore Q\\
P \eif Q, P & \therefore Q
\end{align*}
Claramente, são argumentos válidos. Você pode confirmar que são válidos construindo tabelas-verdade de quatro linhas, mas podemos dizer que eles empregam diferentes \emph{formas} de raciocínio. Seria interessante acompanhar essas diferentes formas de inferência.

Um objetivo de um \emph{sistema de dedução natural} é mostrar que determinados argumentos são válidos, de uma maneira que nos permita compreender o raciocínio que eles podem envolver. Começamos com regras de inferência muito básicas. Essas regras podem ser combinadas para produzir argumentos mais complicados. De fato, esperamos captar todos os argumentos válidos com apenas um pequeno conjunto inicial de regras de inferência.

\emph{Esta é uma maneira muito diferente de pensar sobre argumentos.}

Com tabelas-verdade, consideramos diretamente diferentes maneiras de tornar as sentenças verdadeiras ou falsas. Com sistemas de dedução natural, manipulamos sentenças de acordo com regras que estabelecemos como boas regras. Este último método promete nos dar uma compreensão melhor---ou, pelo menos, uma compreensão diferente---de como os argumentos funcionam.

A adoção da dedução natural pode ser motivada por mais que a busca por compreensão. Ela também pode ser motivada pela \emph{necessidade}. Considere:
\ifHTMLtarget
\begin{earg}
	\item $A_1 \eif C_1$
	\item[\texttherefore] $(A_1 \eand (A_2 \eand (A_3 \eand (A_4 \eand A_5)))) \eif 
	((((C_1 \eor C_2) \eor C_3) \eor C_4) \eor C_5)$
\end{earg}
\else
\begin{align*}
	& A_1 \eif C_1 \\
	\therefore\, & (A_1 \eand (A_2 \eand (A_3 \eand (A_4 \eand A_5)))) \eif {}\\ 
	& \qquad ((((C_1 \eor C_2) \eor C_3) \eor C_4) \eor C_5)
\end{align*}
\fi
Para testar a validade desse argumento, você poderia usar uma tabela-verdade de 1{,}024 linhas. Se fizer isso corretamente, verá que não há linha na qual todas as premissas sejam verdadeiras e a conclusão seja falsa. Assim, saberá que o argumento é válido. (Mas, como acabamos de mencionar, há um sentido em que você não saberá \emph{por que} o argumento é válido.) Considere agora:
\ifHTMLtarget
\begin{earg}
\item $A_1 \eif C_1$
\item[\texttherefore] $(A_1 \eand (A_2 \eand (A_3 \eand (A_4 \eand (A_5 \eand 
 (A_6 \eand (A_7 \eand (A_8 \eand (A_9 \eand A_{10}))))))))) \eif 
 (((((((((C_1 \eor C_2) \eor C_3) \eor C_4) \eor C_5) \eor  C_6) \eor
 C_7) \eor C_8) \eor C_9) \eor C_{10})$
\end{earg}
\else
\begin{align*}
& A_1 \eif C_1 \\
\therefore\, & \small (A_1 \eand (A_2 \eand (A_3 \eand (A_4 \eand (A_5 \eand {}\\
&\qquad (A_6 \eand (A_7 \eand (A_8 \eand (A_9 \eand A_{10}))))))))) \eif {}\\
&(((((((((C_1 \eor C_2) \eor C_3) \eor C_4) \eor C_5) \eor {}\\
& \qquad C_6) \eor C_7) \eor C_8) \eor C_9) \eor C_{10})
\end{align*}
\fi
Esse argumento também é válido---como provavelmente dá para perceber---, mas, para testá-lo, é necessária uma tabela-verdade com $2^{20} = 1{,}048{,}576$ linhas. Em princípio, podemos programar uma máquina para processar tabelas-verdade e informar quando terminar. Na prática, argumentos complicados na LVF podem se tornar \emph{intratáveis} se usarmos tabelas-verdade.

Quando chegarmos à lógica de primeira ordem (LPO) (a partir de \cref{s:FOLBuildingBlocks}), porém, o problema ficará dramaticamente pior. Não há nada semelhante ao teste das tabelas-verdade para a LPO. Para avaliar se um argumento é válido ou não, temos de raciocinar sobre \emph{todas} as interpretações, mas, como veremos, há infinitas interpretações possíveis. Nem em princípio podemos programar uma máquina para processar infinitas interpretações possíveis e informar quando terminar: ela \emph{nunca} terminará. Precisamos ou encontrar uma maneira mais eficiente de raciocinar sobre todas as interpretações, ou procurar algo diferente.

De fato, existem sistemas que codificam maneiras de raciocinar sobre todas as interpretações possíveis. Eles foram desenvolvidos na década de 1950 por Evert Beth e Jaakko Hintikka, mas não seguiremos esse caminho. Em vez disso, recorreremos à dedução natural.

Em vez de raciocinar diretamente sobre todas as valorações (no caso da LVF), tentaremos selecionar algumas regras básicas de inferência. Algumas delas regerão o comportamento dos conectivos sentenciais. Outras regerão o comportamento dos quantificadores e da identidade, que são marcas distintivas da LPO. O sistema de regras resultante nos dará uma nova maneira de pensar sobre a validade de argumentos.
O desenvolvimento moderno da dedução natural remonta a artigos simultâneos e independentes de Gerhard Gentzen e Stanis\l{}aw Ja\'{s}kowski (1934). No entanto, o sistema de dedução natural que consideraremos baseia-se em grande medida no trabalho de Frederic Fitch (publicado pela primeira vez em 1952).


\chapter{Regras básicas para a LVF}\label{s:BasicTFL}
Desenvolveremos um sistema de \define{dedu\c{c}\~{a}o natural}. Para cada conectivo, haverá regras de \define{introdu\c{c}\~{a}o}, que nos permitirão provar uma sentença que tenha esse conectivo como operador lógico principal, e regras de \define{elimina\c{c}\~{a}o}, que nos permitirão provar algo a partir de uma sentença que tenha esse conectivo como operador lógico principal.

\section{A ideia de uma prova formal}
A \emph{prova formal} ou \emph{derivação} é uma sequência de sentenças,
algumas das quais são marcadas como suposições iniciais (ou premissas).
A última linha da prova formal é a conclusão. (Daqui em diante,
chamaremos simplesmente essas sequências de \textquotedblleft provas\textquotedblright{} ou \textquotedblleft derivações\textquotedblright{}, mas tenha em mente que também existem \emph{provas informais}.)

Como ilustração, considere:
	$$\enot (A \eor B) \therefore \enot A \eand \enot B$$
Começaremos uma prova escrevendo a premissa:
\begin{fitchproof}
	\hypo{a1}{\enot (A \eor B)}\PR
\end{fitchproof}
Observe que numeramos a premissa, pois desejaremos voltar a ela. Com efeito, cada linha da prova é numerada, para que possamos fazer referência a ela.

Observe também que traçamos uma linha abaixo da premissa. Tudo o que estiver escrito acima da linha é uma \emph{suposição}. Tudo o que estiver escrito abaixo da linha será algo que decorre das suposições ou será uma nova suposição. Esperamos concluir \textquotedblleft$\enot A \eand \enot B$\textquotedblright{}; portanto, esperamos finalmente concluir nossa prova com
\begin{fitchproof}
	\have[n]{con}{\enot A \eand \enot B}
\end{fitchproof}
para algum número $n$. Não importa em que número de linha terminaremos, mas obviamente preferimos uma prova curta a uma longa.

Suponha, de modo semelhante, que quiséssemos considerar:
$$A\eor B, \enot (A\eand C), \enot (B \eand \enot D) \therefore \enot C\eor D$$
O argumento tem três premissas; portanto, começamos escrevendo todas elas, numeradas, e traçando uma linha abaixo delas:
\begin{fitchproof}
	\hypo{a1}{A \eor B}\PR
	\hypo{a2}{\enot (A\eand C)}\PR
	\hypo{a3}{\enot (B \eand \enot D)}\PR
\end{fitchproof}
e esperamos concluir com alguma linha:
\begin{fitchproof}
	\have[n]{con}{\enot C \eor D}
\end{fitchproof}
Resta apenas explicar cada uma das regras que podemos usar no caminho das premissas à conclusão. As regras são organizadas de acordo com nossos conectivos lógicos.

\section{Reiteração}
A primeira regra é tão estonteantemente óbvia que chega a ser surpreendente que nos demos ao trabalho de apresentá-la.

Se você já mostrou algo no curso de uma prova, a \emph{regra de reiteração} permite repeti-lo em uma nova linha. Por exemplo:
\begin{fitchproof}
	\have[4]{a1}{A \eand B}
	\ellipsesline
	\have[10]{a2}{A \eand B} \by{R}{a1}
\end{fitchproof}
Isso indica que escrevemos \textquotedblleft$A \eand B$\textquotedblright{} na linha~$4$. Agora, em alguma linha posterior — por exemplo, a linha~$10$ — decidimos que queremos repetir isso. Então, escrevemo-lo novamente. Também acrescentamos uma citação que justifica o que escrevemos. Neste caso, escrevemos \textquotedblleft$R$\textquotedblright{}, para indicar que estamos usando a regra de reiteração, e escrevemos \textquotedblleft$4$\textquotedblright{}, para indicar que a aplicamos à linha $4$.

Eis uma expressão geral da regra:
\factoidbox{
\begin{fitchproof}
	\have[m]{a}{\metav{A}}
	\have[\ ]{c}{\metav{A}} \by{R}{a}
\end{fitchproof}}
O ponto é que, se qualquer sentença $\metav{A}$ ocorre em alguma linha,
podemos repetir $\metav{A}$ em linhas posteriores. Cada linha de nossa prova
deve ter uma justificativa, e aqui temos~\textquotedblleft$R$ $m$\textquotedblright{}. Isso significa:
Reiteração, aplicada à linha~$m$.

Duas coisas precisam ser enfatizadas. Primeiro, \textquotedblleft$\metav{A}$\textquotedblright{} não é uma sentença da
LVF. Trata-se, na verdade, de um símbolo da metalinguagem, que usamos quando queremos falar sobre qualquer sentença da LVF (veja \cref{s:UseMention}).
Segundo, e de modo semelhante, \textquotedblleft$m$\textquotedblright{} não é um símbolo que aparecerá em uma
prova. Trata-se, na verdade, de um símbolo da metalinguagem, que usamos quando queremos falar sobre qualquer número de linha de uma prova. Em uma prova real,
as linhas são numeradas \textquotedblleft$1$\textquotedblright{}, \textquotedblleft$2$\textquotedblright{}, \textquotedblleft$3$\textquotedblright{} e assim por diante. Mas, quando
definimos a regra, usamos variáveis como \textquotedblleft$m$\textquotedblright{} para enfatizar que a regra pode ser aplicada em qualquer ponto.

\section{Conjunção}
Suponha que queiramos mostrar que Ludwig é tanto reacionário quanto libertário. Uma maneira óbvia de fazer isso seria a seguinte: primeiro mostramos que Ludwig é reacionário; depois mostramos que Ludwig é libertário; por fim, reunimos essas duas demonstrações para obter a conjunção.

Nosso sistema de dedução natural captará essa ideia de maneira direta. No exemplo dado, podemos adotar a seguinte chave de simbolização:
	\begin{ekey}
		\item[R] Ludwig é reacionário
		\item[L] Ludwig é libertário
	\end{ekey}
Talvez estejamos desenvolvendo uma prova e tenhamos obtido \textquotedblleft$R$\textquotedblright{} na
linha~$8$ e \textquotedblleft$L$\textquotedblright{} na linha~$15$. Então, em qualquer linha subsequente, podemos
obter \textquotedblleft$R \eand L$\textquotedblright{} da seguinte maneira:
\begin{fitchproof}
	\have[8]{a}{R}
	\have[15]{b}{L}
	\have[\ ]{c}{R \eand L} \ai{a, b}
\end{fitchproof}
Observe que cada linha de nossa prova deve ser uma suposição ou
ser justificada por alguma regra. Aqui citamos \textquotedblleft$\eand$I $8$, $15$\textquotedblright{} para
indicar que a linha é obtida pela regra de introdução da conjunção
($\eand$I), aplicada às linhas $8$ e~$15$. Poderíamos igualmente obter:
\begin{fitchproof}
	\have[8]{a}{R}
	\have[15]{b}{L}
	\have[\ ]{c}{L \eand R} \ai{b, a}
\end{fitchproof}
com a citação invertida, para refletir a ordem dos conjuntos. De modo mais geral, eis nossa regra de introdução da conjunção:
\factoidbox{
\begin{fitchproof}
	\have[m]{a}{\metav{A}}
	\have[n]{b}{\metav{B}}
	\have[\ ]{c}{\metav{A}\eand\metav{B}} \ai{a, b}
\end{fitchproof}}

Novamente, este enunciado da regra é \emph{esquemático}. Ele não é
uma prova propriamente dita. \textquotedblleft$\metav{A}$\textquotedblright{} e \textquotedblleft$\metav{B}$\textquotedblright{} são metavariáveis
que representam qualquer sentença da LVF. Do mesmo modo, \textquotedblleft$m$\textquotedblright{} e \textquotedblleft$n$\textquotedblright{} não são
letras que aparecerão em nenhuma prova real, mas símbolos que usamos
para falar sobre qualquer número de linha de qualquer prova. Usamos essas metavariáveis para
enfatizar que a regra pode ser aplicada em qualquer ponto. A regra exige
apenas que ambos os conjuntos estejam disponíveis em algum ponto anterior da
prova. Eles podem estar separados e podem aparecer em qualquer ordem.

A regra é chamada de \textquotedblleft\emph{introdução} da conjunção\textquotedblright{} porque introduz o símbolo \textquotedblleft$\eand$\textquotedblright{} em nossa prova, onde ele talvez estivesse ausente. Correspondentemente, temos uma regra que \emph{elimina} esse símbolo. Suponha que você tenha mostrado que Ludwig é tanto reacionário quanto libertário. Você pode concluir que Ludwig é reacionário. Da mesma forma, pode concluir que Ludwig é libertário. Reunindo isso, obtemos nossas regras de eliminação da conjunção:
\factoidbox{
\begin{fitchproof}
	\have[m]{ab}{\metav{A}\eand\metav{B}}
	\have[\ ]{a}{\metav{A}} \ae{ab}
\end{fitchproof}}
e, de modo equivalente:
\factoidbox{
\begin{fitchproof}
	\have[m]{ab}{\metav{A}\eand\metav{B}}
	\have[\ ]{b}{\metav{B}} \ae{ab}
\end{fitchproof}}
O ponto é simplesmente que, quando você tem uma conjunção em alguma linha de uma prova, pode obter qualquer um dos conjuntos por meio de~{\eand}E. Vale enfatizar um ponto: você só pode aplicar essa regra quando a conjunção é o operador lógico principal. Portanto, não pode inferir \textquotedblleft$D$\textquotedblright{} apenas de \textquotedblleft$C \eor (D \eand E)$\textquotedblright{}!

Mesmo com apenas essas duas regras, já podemos começar a perceber parte do poder de nosso sistema de provas formais. Considere:
\begin{earg}
\item[] $[(A\eor B)\eif(C\eor D)] \eand [(E \eor F) \eif (G\eor H)]$
\item[\texttherefore] $[(E \eor F) \eif (G\eor H)] \eand [(A\eor B)\eif(C\eor D)]$
\end{earg}
O operador lógico principal tanto da premissa quanto da conclusão desse argumento é \textquotedblleft$\eand$\textquotedblright{}. Para fornecer uma prova, começamos escrevendo a premissa, que é nossa suposição. Traçamos uma linha abaixo dela: tudo o que vier depois dessa linha deve decorrer de nossas suposições por meio de (aplicações repetidas de) nossas regras de inferência. Assim, o início da prova é:
\begin{fitchproof}
	\hypo{ab}{{[}(A\eor B)\eif(C\eor D){]} \eand {[}(E \eor F) \eif (G\eor H){]}}\PR
\end{fitchproof}
Da premissa, podemos obter cada um dos conjuntos por meio de {\eand}E. A prova agora é:
\begin{fitchproof}
	\hypo{ab}{{[}(A\eor B)\eif(C\eor D){]} \eand {[}(E \eor F) \eif (G\eor H){]}}\PR
	\have{a}{{[}(A\eor B)\eif(C\eor D){]}} \ae{ab}
	\have{b}{{[}(E \eor F) \eif (G\eor H){]}} \ae{ab}
\end{fitchproof}
Assim, aplicando a regra {\eand}I às linhas 3 e 2 (nessa ordem), chegamos à conclusão desejada. A prova completa é:
\begin{fitchproof}
	\hypo{ab}{{[}(A\eor B)\eif(C\eor D){]} \eand {[}(E \eor F) \eif (G\eor H){]}}\PR

	\have{a}{{[}(A\eor B)\eif(C\eor D){]}} \ae{ab}
	\have{b}{{[}(E \eor F) \eif (G\eor H){]}} \ae{ab}
	\have{ba}{{[}(E \eor F) \eif (G\eor H){]} \eand {[}(A\eor B)\eif(C\eor D){]}} \ai{b,a}
\end{fitchproof}
Esta é uma prova muito simples, mas mostra como podemos encadear regras de prova para formar provas mais longas. A propósito, observe que investigar esse argumento com uma tabela-verdade teria exigido 256 linhas; nossa prova formal exigiu apenas quatro linhas.

Vale a pena dar outro exemplo. Em \cref{s:MoreBracketingConventions}, observamos que este argumento é válido:
	$$A \eand (B \eand C) \therefore (A \eand B) \eand C$$
Para fornecer uma prova correspondente a esse argumento, começamos escrevendo:
\begin{fitchproof}
	\hypo{ab}{A \eand (B \eand C)}\PR
\end{fitchproof}
Da premissa, podemos obter cada um dos conjuntos aplicando~$\eand$E duas vezes. Podemos então aplicar $\eand$E mais duas vezes; assim, nossa prova fica:
\begin{fitchproof}
	\hypo{ab}{A \eand (B \eand C)}\PR
	\have{a}{A} \ae{ab}
	\have{bc}{B \eand C} \ae{ab}
	\have{b}{B} \ae{bc}
	\have{c}{C} \ae{bc}
\end{fitchproof}
Mas agora podemos reintroduzir sem dificuldade as conjunções na ordem que desejávamos; assim, nossa prova final é:
\begin{fitchproof}
	\hypo{abc}{A \eand (B \eand C)}\PR
	\have{a}{A} \ae{abc}
	\have{bc}{B \eand C} \ae{abc}
	\have{b}{B} \ae{bc}
	\have{c}{C} \ae{bc}
	\have{ab}{A \eand B}\ai{a, b}
	\have{con}{(A \eand B) \eand C}\ai{ab, c}
\end{fitchproof}
Lembre-se de que nossa definição oficial de sentenças da LVF permitia apenas conjunções com dois conjuntos. A prova apresentada sugere que poderíamos eliminar os parênteses internos em todas as nossas provas. Entretanto, isso não é padrão e não faremos isso. Em vez disso, manteremos nossas convenções mais austeras para o uso de parênteses. (Ainda assim, continuaremos nos permitindo eliminar os parênteses mais externos, para facilitar a leitura.)

Vejamos uma última ilustração. Ao usar a regra $\eand$I, não há exigência de aplicá-la a sentenças diferentes. Portanto, se quisermos, podemos provar formalmente \textquotedblleft$A \eand A$\textquotedblright{} a partir de \textquotedblleft$A$\textquotedblright{} da seguinte maneira:
\begin{fitchproof}
	\hypo{a}{A}\PR
	\have{aa}{A \eand A}\ai{a, a}
\end{fitchproof}
Simples, mas eficaz.

\section{Condicional}
Considere o seguinte argumento:
\begin{earg}
		\item[] Se Jane é inteligente, então é rápida.
		\item[] Jane é inteligente.
		\item[\texttherefore] Jane é rápida.
\end{earg}
Esse argumento é certamente válido e sugere uma regra direta de eliminação do condicional ($\eif$E):
\factoidbox{
\begin{fitchproof}
	\have[m]{ab}{\metav{A}\eif\metav{B}}
	\have[n]{a}{\metav{A}}
	\have[\ ]{b}{\metav{B}} \ce{ab,a}
\end{fitchproof}}
Essa regra também é chamada às vezes de \emph{modus ponens}. Novamente, trata-se
de uma regra de eliminação, pois permite obter uma sentença que
talvez não contenha \textquotedblleft$\eif$\textquotedblright{}, tendo começado com uma sentença que
continha \textquotedblleft$\eif$\textquotedblright{}. Observe que o condicional $\metav{A}\eif\metav{B}$
e o antecedente~$\metav{A}$ podem estar separados na prova
e podem aparecer em qualquer ordem. É convencional
citar primeiro o número da linha que contém o condicional
$\metav{A}\eif\metav{B}$.

A regra de introdução do condicional também é bastante fácil de motivar. O argumento a seguir deveria ser válido:
	\begin{earg}
		\item Ludwig é reacionário.
		\item[\texttherefore] Se Ludwig é libertário, então Ludwig é reacionário \emph{e} libertário.
	\end{earg}
Se alguém duvidasse da validade desse argumento, poderíamos tentar convencê-lo do contrário explicando nosso raciocínio da seguinte maneira:
	\begin{quote}
		Suponha que Ludwig seja reacionário. Agora, suponha \emph{adicionalmente} que Ludwig seja libertário. Então, pela introdução da conjunção — que acabamos de discutir —, Ludwig é reacionário e libertário. É claro que isso depende da suposição de que Ludwig seja libertário. Mas isso significa simplesmente que, se Ludwig é libertário, então ele é reacionário e libertário.
	\end{quote}
Transpondo para o formato da dedução natural, eis o padrão de raciocínio que acabamos de usar. Começamos com uma premissa, \textquotedblleft{}Ludwig é reacionário\textquotedblright{}, da seguinte maneira:
	\begin{fitchproof}
		\hypo{r}{R}\PR
	\end{fitchproof}
A coisa seguinte que fizemos foi fazer uma suposição \emph{adicional} (\textquotedblleft{}Ludwig é libertário\textquotedblright{}), para fins do argumento. Para indicar que já não estamos lidando \emph{meramente} com nossa premissa original (\textquotedblleft$R$\textquotedblright{}), mas com uma suposição adicional, continuamos nossa prova da seguinte maneira:
	\begin{fitchproof}
		\hypo{r}{R}\PR
		\open
			\hypo{l}{L}\AS
	\end{fitchproof}
Observe que \emph{não} estamos afirmando, na linha~$2$, ter provado
\textquotedblleft$L$\textquotedblright{} a partir da linha~$1$; por isso, não escrevemos nenhuma justificativa para a
suposição adicional na linha~$2$. Contudo, precisamos marcar que ela
é uma suposição adicional. Fazemos isso traçando uma linha abaixo dela
(para indicar que é uma suposição) e recuando-a com uma linha
vertical adicional (para indicar que é adicional).

Com essa suposição extra estabelecida, estamos em condições de
usar~$\eand$I. Assim, podemos continuar nossa prova:
	\begin{fitchproof}
		\hypo{r}{R}\PR
		\open
			\hypo{l}{L}\AS
			\have{rl}{R \eand L}\ai{r, l}
%			\close
%		\have{con}{L \eif (R \eand L)}\ci{l-rl}
	\end{fitchproof}
Assim, mostramos que, sob a suposição adicional \textquotedblleft$L$\textquotedblright{}, podemos obter \textquotedblleft$R \eand L$\textquotedblright{}. Portanto, podemos concluir que, se \textquotedblleft$L$\textquotedblright{} for verdadeira, \textquotedblleft$R \eand L$\textquotedblright{} também será. Ou, de forma mais breve, podemos concluir \textquotedblleft$L \eif (R \eand L)$\textquotedblright{}:
	\begin{fitchproof}
		\hypo{r}{R}\PR
		\open
			\hypo{l}{L}\AS
			\have{rl}{R \eand L}\ai{r, l}
			\close
		\have{con}{L \eif (R \eand L)}\ci{l-rl}
	\end{fitchproof}
Observe que voltamos a usar uma única linha vertical à esquerda. \emph{Descarregamos} a suposição adicional, \textquotedblleft$L$\textquotedblright{}, pois o próprio condicional decorre apenas de nossa suposição original,~\textquotedblleft$R$\textquotedblright{}.

O padrão geral em ação aqui é o seguinte. Primeiro fazemos uma suposição adicional, $\metav{A}$; a partir dessa suposição adicional, provamos~$\metav{B}$. Nesse caso, sabemos o seguinte: se~$\metav{A}$ é verdadeira, então~$\metav{B}$ é verdadeira. Isso é resumido na regra de introdução do condicional:
\factoidbox{
	\begin{fitchproof}
		\open
			\hypo[i]{a}{\metav{A}}\AS
			\have[j]{b}{\metav{B}}
		\close
		\have[\ ]{ab}{\metav{A}\eif\metav{B}}\ci{a-b}
\end{fitchproof}}
Pode haver quantas linhas você quiser — ou tão poucas quanto quiser — entre as linhas $i$
e~$j$.

Será útil oferecer uma segunda ilustração da aplicação de $\eif$I. Suponha que queiramos considerar o seguinte:
	$$P \eif Q, Q \eif R \therefore P \eif R$$
Começamos listando \emph{ambas} as premissas. Então, como queremos chegar a um condicional (a saber, \textquotedblleft$P \eif R$\textquotedblright{}), fazemos adicionalmente a suposição do antecedente desse condicional. Assim, nossa prova principal começa:
\begin{fitchproof}
	\hypo{pq}{P \eif Q}\PR
	\hypo{qr}{Q \eif R}\PR
	\open
		\hypo{p}{P}\AS
	\close
\end{fitchproof}
Observe que tornamos \textquotedblleft$P$\textquotedblright{} disponível tratando-o como uma suposição adicional; agora, podemos usar {\eif}E na primeira premissa. Isso produzirá \textquotedblleft$Q$\textquotedblright{}. Podemos então usar {\eif}E na segunda premissa. Assim, ao supor \textquotedblleft$P$\textquotedblright{}, conseguimos provar \textquotedblleft$R$\textquotedblright{}; portanto, aplicamos a regra {\eif}I — descarregando \textquotedblleft$P$\textquotedblright{} — e concluímos a prova. Reunindo tudo isso, temos:
\label{HSproof}
\begin{fitchproof}
	\hypo{pq}{P \eif Q}\PR
	\hypo{qr}{Q \eif R}\PR
	\open
		\hypo{p}{P}\AS
		\have{q}{Q}\ce{pq,p}
		\have{r}{R}\ce{qr,q}
	\close
	\have{pr}{P \eif R}\ci{p-r}
\end{fitchproof}


\section{Suposições adicionais e subprovas}
A regra $\eif$I evocou a ideia de fazer suposições adicionais. Elas precisam ser tratadas com algum cuidado. Considere esta prova:
\begin{fitchproof}
	\hypo{a}{A}\PR
	\open
		\hypo{b1}{B}\AS
		\have{b2}{B} \by{R}{b1}
	\close
	\have{con}{B \eif B}\ci{b1-b2}
\end{fitchproof}
Isso está perfeitamente de acordo com as regras que estabelecemos até aqui e não deve parecer particularmente estranho. Como \textquotedblleft$B \eif B$\textquotedblright{} é uma tautologia, nenhuma premissa específica deveria ser necessária para prová-la.

Mas suponha que agora tentássemos continuar a prova da seguinte maneira:
\begin{fitchproof}
	\hypo{a}{A}\PR
	\open
		\hypo{b1}{B}\AS
		\have{b2}{B} \by{R}{b1}
	\close
	\have{con}{B \eif B}\ci{b1-b2}
\ifHTMLtarget
	\have{b}{B} \by{tentativa indevida de aplicar $\eif$E}{con, b2}
\else
	\have{b}{B} \by{tentativa indevida}{}
	\have [\ ]{x}{} \by{de aplicar $\eif$E}{con, b2}
\fi
\end{fitchproof}
Se tivéssemos permissão para fazer isso, seria um desastre. Isso nos permitiria provar qualquer letra sentencial a partir de qualquer outra letra sentencial. Entretanto, se você me disser que Anne é rápida (simbolizada por \textquotedblleft$A$\textquotedblright{}), não deveríamos poder concluir que a rainha Boudica tinha vinte pés de altura (simbolizada por \textquotedblleft$B$\textquotedblright{})! Precisamos proibir isso, mas como implementar a proibição?

Podemos descrever o processo de fazer uma suposição adicional como a realização de uma \emph{subprova}: uma prova subsidiária dentro da prova principal. Quando iniciamos uma subprova, traçamos outra linha vertical para indicar que já não estamos na prova principal. Em seguida, escrevemos a suposição na qual a subprova se baseará. Uma subprova pode ser entendida essencialmente como a formulação da seguinte pergunta: \emph{o que poderíamos mostrar se também fizéssemos essa suposição adicional?}

Quando estamos trabalhando dentro da subprova, podemos fazer referência à suposição adicional que fizemos ao introduzir a subprova e a qualquer coisa que obtivemos a partir de nossas suposições originais. (Afinal, essas suposições originais continuam em vigor.) Contudo, em algum momento, queremos deixar de trabalhar com a suposição adicional: queremos retornar da subprova à prova principal. Para indicar que retornamos à prova principal, a linha vertical da subprova chega ao fim. Nesse ponto, dizemos que a subprova está \define{fechada}. Ao fechar uma subprova, deixamos de lado a suposição adicional; portanto, será ilegítimo recorrer a qualquer coisa que dependa dessa suposição adicional. Assim, estipulamos:
\factoidbox{Para citar uma linha individual ao aplicar uma regra:
\begin{compactlist}
\item a linha deve vir antes da linha em que a regra é aplicada, mas
\item não pode ocorrer dentro de uma subprova que tenha sido fechada antes da linha em que a regra é aplicada.
\end{compactlist}}
Essa estipulação exclui a tentativa de prova desastrosa apresentada acima. A aplicação da regra $\eif$E na linha~$5$ exige que citemos duas linhas individuais anteriores na prova. Uma dessas linhas (a saber, a linha~$3$) ocorre dentro de uma subprova (linhas $2$--$3$). Na linha $5$, quando temos de citar a linha $3$, a subprova já foi fechada. Isso é ilegítimo: não temos permissão para citar a linha $3$ na linha $5$.

O fechamento de uma subprova é chamado de \define{descarga} das suposições
dessa subprova. Assim, podemos formular o ponto desta maneira: \emph{você não pode
fazer referência a nada que tenha sido obtido usando uma suposição descarregada}.

As subprovas, portanto, permitem pensar sobre o que poderíamos mostrar se fizéssemos suposições adicionais. A lição a extrair disso não é surpreendente — no curso de uma prova, precisamos acompanhar com muito cuidado quais suposições estamos fazendo a cada momento. Nosso sistema de provas faz isso de maneira muito gráfica. (De fato, é precisamente por isso que escolhemos usar \emph{este} sistema de provas.)

Uma vez que começamos a pensar sobre o que podemos mostrar fazendo suposições adicionais, nada nos impede de formular a pergunta: o que poderíamos mostrar se fizéssemos \emph{ainda mais} suposições? Isso pode nos motivar a introduzir uma subprova dentro de outra subprova. Eis um exemplo que usa apenas as regras consideradas até aqui:
\begin{fitchproof}
\hypo{a}{A}\PR
\open
	\hypo{b}{B}\AS
	\open
		\hypo{c}{C}\AS
		\have{ab}{A \eand B}\ai{a,b}
	\close
	\have{cab}{C \eif (A \eand B)}\ci{c-ab}
\close
\have{bcab}{B \eif (C \eif (A \eand B))}\ci{b-cab}
\end{fitchproof}
Observe que a citação na linha~$4$ faz referência à suposição inicial
(na linha~$1$) e a uma suposição de uma subprova (na
linha~$2$). Isso está perfeitamente correto, pois nenhuma das duas suposições havia
sido descarregada naquele momento (isto é, até a linha~$4$).

Ainda assim, precisamos acompanhar cuidadosamente o que estamos supondo a cada momento. Suponha que tentássemos continuar a prova da seguinte maneira:
\begin{fitchproof}
\hypo{a}{A}\PR
\open
	\hypo{b}{B}\AS
	\open
		\hypo{c}{C}\AS
		\have{ab}{A \eand B}\ai{a,b}
	\close
	\have{cab}{C \eif (A \eand B)}\ci{c-ab}
\close
\have{bcab}{B \eif(C \eif (A \eand B))}\ci{b-cab}
\ifHTMLtarget
	\have{bcab}{C \eif (A \eand B)}\by{tentativa indevida de aplicar $\eif$I}{c-ab}
\else
	\have{bcab}{C \eif (A \eand B)}\by{tentativa indevida}{}
	\have [\ ]{x}{} \by{de aplicar $\eif$I}{c-ab}
\fi
\end{fitchproof}
Isso seria terrível. Se lhe dissermos que Anne é inteligente, você não
deveria poder inferir que, se Cath é inteligente (simbolizada por \textquotedblleft$C$\textquotedblright{}), então
\emph{tanto} Anne é inteligente quanto a rainha Boudica tinha 20~pés de altura! Mas
é exatamente isso que uma prova desse tipo sugeriria, se ela fosse permitida.

O problema essencial é que a subprova que começou com a suposição~\textquotedblleft$C$\textquotedblright{} dependia crucialmente do fato de termos suposto \textquotedblleft$B$\textquotedblright{} na linha~$2$. Na linha~$6$, \emph{descarregamos} a suposição~\textquotedblleft$B$\textquotedblright{}: deixamos de perguntar o que poderíamos mostrar se também supuséssemos \textquotedblleft$B$\textquotedblright{}. Portanto, é simplesmente trapacear tentar nos valer (na linha~$7$) da subprova que começou com a suposição~\textquotedblleft$C$\textquotedblright{}. Assim, estipulamos, como antes, que uma subprova só pode ser citada em uma linha se não estiver contida em outra subprova que já esteja fechada nessa linha. A tentativa de prova desastrosa viola essa estipulação. A subprova das linhas $3$--$4$ ocorre dentro de uma subprova que termina na linha~$5$. Portanto, ela não pode ser invocada na linha~$7$.

Eis mais um caso que precisamos excluir:
\begin{fitchproof}
\hypo{a}{A}\PR
\open
	\hypo{b}{B}\AS
	\open
	\hypo{c}{C}\AS
	\have{bc}{B \eand C}\ai{b,c}
	\have{c2}{C}\ae{bc}
	\close
\close
\ifHTMLtarget
\have{bcab}{B \eif C}\by{tentativa indevida de aplicar $\eif$I}{b-c2}
\else
\have{bcab}{B \eif C}\by{tentativa indevida}{}
\have [\ ]{x}{} \by{de aplicar $\eif$I}{b-c2}
\fi
\end{fitchproof}
Aqui estamos tentando citar uma subprova que começa na linha~$2$ e termina na linha~$5$ — mas a sentença da linha~$5$ depende não apenas da suposição da linha~$2$, mas também de outra suposição (a linha~$3$) que não descarregamos ao final da subprova. A subprova iniciada na linha~$3$ ainda está aberta na linha~$5$. Entretanto, $\eif$I exige que a última linha da subprova dependa \emph{somente} da suposição da subprova citada, isto é, da subprova que começa na linha~$2$ (e de tudo o que vem antes dela), e não das suposições de quaisquer subprovas contidas nela. Em particular, a última linha da subprova citada não pode estar, ela própria, dentro de uma subprova aninhada.

\factoidbox{Para citar uma subprova ao aplicar uma regra:
\begin{compactlist} 
\item a subprova citada deve vir inteiramente antes da aplicação da regra em que é citada,
\item a subprova citada não deve estar contida em outra subprova que esteja fechada na linha em que é citada, e
\item a última linha da subprova citada não deve ocorrer dentro de uma subprova aninhada.
\end{compactlist}}

Um último ponto para enfatizar como as regras podem ser aplicadas: quando uma regra exige que você cite uma linha individual, não pode citar uma subprova em seu lugar; e, quando exige que você cite uma subprova, não pode citar uma linha individual em seu lugar. Por exemplo, isto está incorreto:
\begin{fitchproof}
\hypo{a}{A}\PR
\open
	\hypo{b}{B}\AS
	\open
	\hypo{c}{C}\AS
	\have{bc}{B \eand C}\ai{b,c}
	\have{c2}{C}\ae{bc}
	\close
\ifHTMLtarget
	\have{c3}{C}\by{tentativa indevida de aplicar R}{c-c2}
\else
\have{c3}{C}\by{tentativa indevida}{}
\have [\ ]{x}{} \by{de aplicar R}{c-c2}
\fi
\close
\have[7]{bcab}{B \eif C}\ci{b-c3}
\end{fitchproof}
Aqui, tentamos justificar $C$ na linha~$6$ pela regra de reiteração,
mas citamos com ela a subprova das linhas $3$--$5$. Essa subprova está
fechada e, em princípio, pode ser citada na linha~$6$. (Por
exemplo, poderíamos usá-la para justificar $C \eif C$ por $\eif$I.) Mas a
regra de reiteração~R exige que você cite uma linha individual; portanto, citar
a subprova inteira é inadmissível (mesmo que essa subprova contenha
a sentença~$C$ que queremos reiterar). Citar as linhas individuais
que contêm $C$ ($3$ ou $5$) na linha~$6$ também é, evidentemente, inadmissível,
pois elas estão em uma subprova fechada antes da linha~$6$.


É sempre permitido abrir uma subprova com qualquer suposição. Entretanto, há alguma estratégia envolvida na escolha de uma suposição útil. Começar uma subprova com uma suposição arbitrária e absurda apenas desperdiçaria linhas da prova. Para obter um condicional por {\eif}I, por exemplo, você deve supor o antecedente do condicional em uma subprova.

Da mesma forma, é sempre permitido fechar uma subprova (e descarregar
suas suposições). Contudo, não será útil fazê-lo até que você tenha
alcançado algo útil. Uma vez fechada a subprova, você só pode citar
a subprova inteira em qualquer justificativa. Essas regras que
exigem uma subprova ou subprovas, por sua vez, exigem que a última linha
da subprova seja uma sentença de alguma forma. Por exemplo, só é
permitido citar uma subprova para~$\eif$I se a linha que você está
justificando tiver a forma $\metav{A} \eif \metav{B}$, $\metav{A}$ for
a suposição de sua subprova e $\metav{B}$ for a última linha de
sua subprova.


\section{Bicondicional}

As regras para o bicondicional serão versões de duas vias das regras
para o condicional.

Para provar, por exemplo, \textquotedblleft$F \eiff G$\textquotedblright{}, você deve ser capaz de
provar \textquotedblleft$G$\textquotedblright{} sob a suposição \textquotedblleft$F$\textquotedblright{} e \emph{também} provar \textquotedblleft$F$\textquotedblright{} sob a
suposição~\textquotedblleft$G$\textquotedblright{}. Portanto, a regra de introdução do bicondicional ({\eiff}I)
exige duas subprovas. Esquematicamente, a regra funciona assim:
\factoidbox{
\begin{fitchproof}
	\open
		\hypo[i]{a1}{\metav{A}}\AS
		\have[j]{b1}{\metav{B}}
	\close
	\open
		\hypo[k]{b2}{\metav{B}}\AS
		\have[l]{a2}{\metav{A}}
	\close
	\have[\ ]{ab}{\metav{A}\eiff\metav{B}}\bi{a1-b1,b2-a2}
\end{fitchproof}}
Pode haver quantas linhas você quiser entre $i$ e $j$, e quantas linhas você quiser entre $k$ e $l$. Além disso, as subprovas podem aparecer em qualquer ordem, e a segunda não precisa vir imediatamente depois da primeira.

A regra de eliminação do bicondicional ({\eiff}E) permite fazer um pouco mais do que a regra do condicional. Se você tem a subsentença do lado esquerdo do bicondicional, pode obter a subsentença do lado direito. Se tem a subsentença do lado direito, pode obter a do lado esquerdo. Portanto, permitimos:
\factoidbox{
\begin{fitchproof}
	\have[m]{ab}{\metav{A}\eiff\metav{B}}
	\have[n]{a}{\metav{A}}
	\have[\ ]{b}{\metav{B}} \be{ab,a}
\end{fitchproof}}
e, de modo equivalente:
\factoidbox{\begin{fitchproof}
	\have[m]{ab}{\metav{A}\eiff\metav{B}}
	\have[n]{a}{\metav{B}}
	\have[\ ]{b}{\metav{A}} \be{ab,a}
\end{fitchproof}}
Observe que o bicondicional e a metade direita ou esquerda podem estar
separados uns dos outros e podem aparecer em qualquer ordem. Contudo,
na citação de~$\eiff$E, é boa prática citar primeiro o bicondicional.

\section{Disjunção}
Suponha que Ludwig seja reacionário. Então Ludwig é reacionário ou libertário. Afinal, dizer que Ludwig é reacionário ou libertário é dizer algo mais fraco do que dizer que Ludwig é reacionário.

Vamos enfatizar esse ponto. Suponha que Ludwig seja reacionário. Segue-se que Ludwig é \emph{ou} reacionário \emph{ou} uma quincã. Do mesmo modo, segue-se que \emph{ou} Ludwig é reacionário \emph{ou} as quincãs são a única fruta. Igualmente, segue-se que \emph{ou} Ludwig é reacionário \emph{ou} Deus está morto. Muitas dessas são inferências estranhas, mas não há nada \emph{logicamente} errado com elas (mesmo que talvez violem todo tipo de norma conversacional implícita).

Munidos de tudo isso, apresentamos a(s) regra(s) de introdução da disjunção:
\factoidbox{\begin{fitchproof}
	\have[m]{a}{\metav{A}}
	\have[\ ]{ab}{\metav{A}\eor\metav{B}}\oi{a}
\end{fitchproof}}
e
\factoidbox{\begin{fitchproof}
	\have[m]{a}{\metav{A}}
	\have[\ ]{ba}{\metav{B}\eor\metav{A}}\oi{a}
\end{fitchproof}}
Observe que \metav{B} pode ser \emph{qualquer} sentença, de modo que a seguinte é uma prova perfeitamente aceitável:
\begin{fitchproof}
	\hypo{m}{M}\PR
	\have{mmm}{M \eor ([(A\eiff B) \eif (C \eand D)] \eiff [E \eand F])}\oi{m}
\end{fitchproof}
Mostrar isso com uma tabela-verdade teria exigido 128 linhas.

A regra de eliminação da disjunção é um pouco mais complicada. Suponha que Ludwig seja reacionário ou libertário. O que você pode concluir? Não que Ludwig seja reacionário; ele pode ser libertário. Tampouco que Ludwig seja libertário; ele pode ser simplesmente reacionário. As disjunções, por si mesmas, são difíceis de manipular.

Mas suponha que pudéssemos mostrar ambas as coisas seguintes: primeiro, que o fato de Ludwig ser reacionário implica que ele é um economista austríaco; segundo, que o fato de Ludwig ser libertário implica que ele é um economista austríaco. Então, se sabemos que Ludwig é reacionário ou libertário, sabemos que, qualquer que seja o caso, Ludwig é um economista austríaco. Essa ideia pode ser expressa na regra seguinte, nossa regra de eliminação da disjunção ($\eor$E):
\factoidbox{
	\begin{fitchproof}
		\have[m]{ab}{\metav{A}\eor\metav{B}}
		\open
			\hypo[i]{a}{\metav{A}} \AS
			\have[j]{c1}{\metav{C}}
		\close
		\open
			\hypo[k]{b}{\metav{B}} \AS
			\have[l]{c2}{\metav{C}}
		\close
		\have[ ]{c}{\metav{C}}\oe{ab, a-c1,b-c2}
	\end{fitchproof}}
É claro que escrever isso é um pouco mais trabalhoso do que escrever nossas regras anteriores, mas o ponto é bastante simples. Suponha que tenhamos uma disjunção, $\metav{A} \eor \metav{B}$. Suponha que tenhamos duas subprovas mostrando que $\metav{C}$ decorre da suposição de~$\metav{A}$ e que $\metav{C}$ decorre da suposição de~$\metav{B}$. Então podemos inferir a própria sentença $\metav{C}$. Como de costume, pode haver quantas linhas você quiser entre $i$ e $j$, e quantas linhas você quiser entre $k$ e $l$. Além disso, as subprovas e a disjunção podem aparecer em qualquer ordem e não precisam ser adjacentes.

Alguns exemplos podem ajudar a ilustrar isso. Considere o argumento:
$$(P \eand Q) \eor (P \eand R) \therefore P$$
Uma prova de exemplo poderia ser:
	\begin{fitchproof}
		\hypo{prem}{(P \eand Q) \eor (P \eand R) }\PR
			\open
				\hypo{pq}{P \eand Q}\AS
				\have{p1}{P}\ae{pq}
			\close
			\open
				\hypo{pr}{P \eand R}\AS
				\have{p2}{P}\ae{pr}
			\close
		\have{con}{P}\oe{prem, pq-p1, pr-p2}
	\end{fitchproof}
Eis um exemplo um pouco mais difícil. Considere:
	$$ A \eand (B \eor C) \therefore (A \eand B) \eor (A \eand C)$$
Eis uma prova correspondente a esse argumento:
	\begin{fitchproof}
		\hypo{aboc}{A \eand (B \eor C)}\PR
		\have{a}{A}\ae{aboc}
		\have{boc}{B \eor C}\ae{aboc}
		\open
			\hypo{b}{B}\AS
			\have{ab}{A \eand B}\ai{a,b}
			\have{abo}{(A \eand B) \eor (A \eand C)}\oi{ab}
		\close
		\open
			\hypo{c}{C}\AS
			\have{ac}{A \eand C}\ai{a,c}
			\have{aco}{(A \eand B) \eor (A \eand C)}\oi{ac}
		\close
	\have{con}{(A \eand B) \eor (A \eand C)}\oe{boc, b-abo, c-aco}
	\end{fitchproof}
Não se assuste se achar que não teria sido capaz de elaborar essa prova por conta própria. A capacidade de elaborar provas novas vem com a prática, e abordaremos algumas estratégias para encontrar provas em \cref{s:stratTFL}. A questão principal, neste momento, é se, olhando para a prova, você consegue ver que ela está de acordo com as regras que estabelecemos. Isso envolve apenas verificar cada linha e certificar-se de que ela é justificada de acordo com essas regras.


\section{Contradição e negação}\label{sec:contradiction}

Resta tratar de apenas um conectivo: a negação. Para abordá-la, porém, precisamos relacionar a negação à \emph{contradição}.

Uma forma eficaz de argumentar é levar seu oponente a se contradizer. Nesse ponto, ele está encurralado. Tem de abandonar pelo menos uma de suas suposições. Usaremos essa ideia em nosso sistema de provas, acrescentando um novo símbolo, \textquotedblleft$\ered$\textquotedblright{}, às nossas provas. Ele deve ser lido como algo semelhante a \textquotedblleft{}contradição!\textquotedblright{}, \textquotedblleft{}reductio!\textquotedblright{} ou \textquotedblleft{}mas isso é absurdo!\textquotedblright{}. A regra para introduzir esse símbolo é que podemos usá-lo sempre que nos contradizemos explicitamente, isto é, sempre que encontramos tanto uma sentença quanto sua negação em nossa prova:
\factoidbox{
\begin{fitchproof}
  \have[m]{na}{\enot\metav{A}}
  \have[n]{a}{\metav{A}}
  \have[ ]{bot}{\ered}\ne{na, a}
\end{fitchproof}}
Não importa em que ordem a sentença e sua negação aparecem, e elas não
precisam estar em linhas adjacentes. É convencional citar primeiro o
número da linha da negação, seguido pelo número da linha da sentença da
qual ela é a negação.

Há, evidentemente, uma estreita relação entre contradição e negação.
A regra $\enot$E permite passar de duas sentenças contraditórias---$\metav{A}$ e sua negação $\enot \metav{A}$---a uma contradição explícita~$\ered$. Escolhemos esse rótulo por uma razão: essa é a regra mais básica que permite passar de uma premissa que contém uma negação, isto é, $\enot\metav{A}$, a uma sentença que não a contém, isto é, $\ered$. Portanto, é uma regra que \emph{elimina}~$\enot$.

Dissemos que \textquotedblleft$\ered$\textquotedblright{} deve ser lido como algo semelhante a \textquotedblleft{}contradição!\textquotedblright{}, mas isso não nos diz muita coisa sobre o símbolo. Há, grosso modo, três maneiras de entendê-lo.
	\begin{compactlist}
		\item Podemos considerar \textquotedblleft$\ered$\textquotedblright{} uma nova sentença atômica da LVF, mas uma que só pode ter o valor de verdade Falso.
		\item Podemos considerar \textquotedblleft$\ered$\textquotedblright{} uma abreviação para alguma contradição canônica, como \textquotedblleft$A \eand \enot A$\textquotedblright{}. Isso terá o mesmo efeito que a opção anterior---evidentemente, \textquotedblleft$A \eand \enot A$\textquotedblright{} só pode ter o valor de verdade Falso---, mas significa que, oficialmente, não precisamos acrescentar um novo símbolo à LVF.
		\item Podemos considerar \textquotedblleft$\ered$\textquotedblright{}, não um símbolo da LVF, mas algo mais parecido com um \emph{sinal de pontuação} que aparece em nossas provas. (Ele estaria no mesmo nível dos números de linha e das linhas verticais, por assim dizer.)
	\end{compactlist}
Há algo filosoficamente muito atraente na terceira opção, mas aqui adotaremos \emph{oficialmente} a primeira. \textquotedblleft$\ered$\textquotedblright{} deve ser lido como uma letra sentencial que é sempre falsa. Isso significa que podemos manipulá-lo, em nossas provas, como qualquer outra sentença.

Ainda precisamos enunciar uma regra para a introdução da negação. A regra é muito simples: se supor algo leva a uma contradição, então a suposição deve estar errada. Essa ideia motiva a regra seguinte:
\factoidbox{\begin{fitchproof}
\open
	\hypo[i]{a}{\metav{A}}\AS
	\have[j]{nb}{\ered}
\close
\have[\ ]{na}{\enot\metav{A}}\ni{a-nb}
\end{fitchproof}}
Pode haver quantas linhas você quiser entre $i$ e $j$. Para ver isso na prática e em interação com a negação, considere a prova:
	\begin{fitchproof}
		\hypo{d}{D}\AS
		\open
			\hypo{nd}{\enot D}\AS
			\have{ndr}{\ered}\ne{nd, d}
		\close
		\have{con}{\enot\enot D}\ni{nd-ndr}
	\end{fitchproof}

Se a suposição de que $\metav{A}$ é verdadeira leva a uma contradição, $\metav{A}$ não pode ser verdadeira; isto é, deve ser falsa; isto é, $\enot\metav{A}$ deve ser verdadeira. Naturalmente, se a suposição de que $\metav{A}$ é falsa (isto é, a suposição de que $\enot\metav{A}$ é verdadeira) leva a uma contradição, então $\metav{A}$ não pode ser falsa; isto é, deve ser verdadeira. Assim, podemos considerar a regra seguinte:
\factoidbox{\begin{fitchproof}
\open
	\hypo[i]{a}{\enot\metav{A}}\AS
	\have[j]{nb}{\ered}
\close
\have[\ ]{na}{\metav{A}}\ip{a-nb}
\end{fitchproof}}
Essa regra é chamada de \emph{prova indireta}, pois permite provar $\metav{A}$ indiretamente, supondo sua negação. Formalmente, a regra é muito semelhante a $\enot$I, mas $\metav{A}$ e $\enot\metav{A}$ trocaram de lugar. Como $\enot\metav{A}$ não é a conclusão da regra, não estamos introduzindo~$\enot$; portanto, IP não é uma regra que introduz algum conectivo. Ela também não elimina um conectivo, pois não tem premissas independentes que contenham~$\enot$, mas apenas uma subprova com uma suposição da forma~$\enot\metav{A}$. Em contraste, $\enot$E tem uma premissa da forma $\enot\metav{A}$: é por isso que $\enot$E elimina~$\enot$, mas IP não.
\footnote{Há lógicos que têm reservas quanto à IP, mas não quanto a $\enot$E. Eles são chamados de \textquotedblleft{}intuicionistas\textquotedblright{}. Os intuicionistas não aceitam nossa suposição básica de que toda sentença tem um de dois valores de verdade, verdadeiro ou falso. Eles também acham que $\enot$ funciona de maneira diferente---para eles, uma prova de $\ered$ a partir de $\metav{A}$ garante $\enot \metav{A}$, mas uma prova de $\ered$ a partir de $\enot\metav{A}$ não garante $\metav{A}$, apenas $\enot\enot\metav{A}$. Portanto, para eles, $\metav{A}$ e $\enot\enot\metav{A}$ não são equivalentes.}

Usando $\enot$I, conseguimos dar uma prova de $\enot\enot\metav{D}$ a partir de $\metav{D}$. Usando IP, podemos seguir na direção oposta (com uma prova essencialmente igual).
	\begin{fitchproof}
		\hypo{d}{\enot\enot D}\PR
		\open
			\hypo{nd}{\enot D}\AS
			\have{ndr}{\ered}\ne{d, nd}
		\close
		\have{con}{D}\ip{nd-ndr}
	\end{fitchproof}

Precisamos de uma última regra. Ela é uma espécie de regra de eliminação para \textquotedblleft$\ered$\textquotedblright{} e é conhecida como \emph{explosão}.\footnote{O nome latino desse princípio é \emph{ex contradictione quodlibet}, \textquotedblleft{}da contradição, qualquer coisa\textquotedblright{}.} Se obtemos uma contradição, simbolizada por \textquotedblleft$\ered$\textquotedblright{}, então podemos inferir o que quisermos. Como isso pode ser motivado como uma regra de argumentação? Considere o recurso retórico \textquotedblleft{}\ldots e, se \emph{isso} for verdade, eu como meu chapéu\textquotedblright{}. Como as contradições simplesmente não podem ser verdadeiras, se uma \emph{for} verdadeira, não apenas comerei meu chapéu, mas também o terei comigo. Eis a regra formal:
\factoidbox{\begin{fitchproof}
\have[m]{bot}{\ered}
\have[ ]{}{\metav{A}}\re{bot}
\end{fitchproof}}
Observe que \metav{A} pode ser \emph{qualquer} sentença.

A regra da explosão é um pouco estranha. Parece que $\metav{A}$ surge em nossa prova como um coelho saindo da cartola. Ao tentar encontrar provas, é muito tentador usá-la em toda parte, pois parece muito poderosa. Resista a essa tentação: você só pode aplicá-la quando já tem~$\ered$! E só obtemos $\ered$ quando nossas suposições são contraditórias.

Ainda assim, não é estranho que, de uma contradição, se siga qualquer coisa? Não de acordo com nossa noção de implicação e validade. $\metav{A}$ implica $\metav{B}$ \ifeff{} não há valoração das letras sentenciais que torne $\metav{A}$ verdadeira e $\metav{B}$ falsa ao mesmo tempo. Ora, $\ered$ é uma contradição---nunca é verdadeira, qualquer que seja a valoração das letras sentenciais. Como não há valoração que torne $\ered$ verdadeira, naturalmente também não há valoração que torne $\ered$ verdadeira e $\metav{B}$ falsa! Assim, de acordo com nossa definição de implicação, $\ered \entails \metav{B}$, qualquer que seja $\metav{B}$. Uma contradição implica qualquer coisa.\footnote{Há lógicos que não aceitam isso. Eles acham que, se $\metav{A}$ implica $\metav{B}$, deve haver alguma \emph{conexão relevante} entre $\metav{A}$ e $\metav{B}$---e não há nenhuma entre $\ered$ e uma sentença arbitrária~$\metav{B}$. Assim, esses lógicos desenvolvem outras lógicas, \textquotedblleft{}relevantes\textquotedblright{}, nas quais a regra da explosão não é permitida.}

\emph{Estas são todas as regras básicas do sistema de provas da LVF.}

\practiceproblems

\problempart
As duas \textquotedblleft{}provas\textquotedblright{} seguintes estão \emph{incorretas}. Explique os erros que elas cometem.
\begin{fitchproof}
\hypo{abc}{(\enot L \eand A) \eor L}\PR
\open
\hypo{nla}{\enot L \eand A}\AS
\have{nl}{\enot L}\ae{nla}
	\have{a}{A}\ae{abc}
\close
\open
	\hypo{l}{L}\AS
	\have{red}{\ered}\ne{nl, l}
	\have{a2}{A}\re{red}
\close
\have{con}{A}\oe{abc, nla-a, l-a2}
\end{fitchproof}

\begin{fitchproof}
\hypo{abc}{A \eand (B \eand C)}\PR
\hypo{bcd}{(B \eor C) \eif D}\PR
\have{b}{B}\ae{abc}
\have{bc}{B \eor C}\oi{b}
\have{d}{D}\ce{bc, bcd}
\end{fitchproof}

\problempart
As três provas seguintes estão sem suas citações (a regra e os números das linhas). Acrescente-as para transformá-las em provas \emph{bona fide}. Além disso, escreva o argumento correspondente a cada prova.
\begin{multicols}{2}
\begin{fitchproof}
\hypo{ps}{P \eand S}
\hypo{nsor}{S \eif R}
\have{p}{P}%\ae{ps}
\have{s}{S}%\ae{ps}
\have{r}{R}%\ce{nsor, s}
\have{re}{R \eor E}%\oi{r}
\end{fitchproof}

\begin{fitchproof}
\hypo{ad}{A \eif D}
\open
	\hypo{ab}{A \eand B}
	\have{a}{A}%\ae{ab}
	\have{d}{D}%\ce{ad, a}
	\have{de}{D \eor E}%\oi{d}
\close
\have{conc}{(A \eand B) \eif (D \eor E)}%\ci{ab-de}
\end{fitchproof}

\begin{fitchproof}
\hypo{nlcjol}{\enot L \eif (J \eor L)}
\hypo{nl}{\enot L}
\have{jol}{J \eor L}%\ce{nlcjol, nl}
\open
	\hypo{j}{J}
	\have{jj}{J \eand J}%\ai{j}
	\have{j2}{J}%\ae{jj}
\close
\open
	\hypo{l}{L}
	\have{red}{\ered}%\ne{nl, l}
	\have{j3}{J}%\re{red}
\close
\have{conc}{J}%\oe{jol, j-j2, l-j3}
\end{fitchproof}
\end{multicols}

\solutions
\problempart
\label{pr.solvedTFLproofs}
Forneça uma prova para cada um dos argumentos a seguir:
\begin{compactlist}
\item $J\eif\enot J \therefore \enot J$
\item $Q\eif(Q\eand\enot Q) \therefore \enot Q$
\item $A\eif (B\eif C) \therefore (A\eand B)\eif C$
\item $K\eand L \therefore K\eiff L$
\item $(C\eand D)\eor E \therefore E\eor D$
\item $A\eiff B, B\eiff C \therefore A\eiff C$
\item $\enot F\eif G, F\eif H \therefore G\eor H$
\item $(Z\eand K) \eor (K\eand M), K \eif D \therefore D$
\item $P \eand (Q\eor R), P\eif \enot R \therefore Q\eor E$
\item $S\eiff T \therefore S\eiff (T\eor S)$
\item $\enot (P \eif Q) \therefore \enot Q$
\item $\enot (P \eif Q) \therefore P$
\end{compactlist}


\chapter{Construindo provas}\label{s:stratTFL}

Não há uma receita simples para encontrar provas, e nada substitui a prática. Ainda assim, seguem algumas regras práticas e estratégias que convém ter em mente.

\section{Trabalhando de trás para frente a partir do que queremos}

Suponha que você esteja tentando encontrar uma prova para alguma conclusão~$\metav{C}$, que será a última linha da sua prova. A primeira coisa a fazer é olhar para~$\metav{C}$ e perguntar qual é a regra de introdução do seu operador lógico principal. Isso dá uma ideia do que deve acontecer \emph{antes} da última linha da prova. As justificativas para a regra de introdução exigem uma ou duas outras sentenças acima da última linha, ou uma ou duas subprovas. Além disso, a partir de~$\metav{C}$ você consegue determinar quais são essas sentenças, ou quais são as suposições e conclusões da(s) subprova(s). Então, você pode escrever essas sentenças ou esquematizar a(s) subprova(s) acima da última linha e tratá-las como seus novos objetivos.

Por exemplo: se a sua conclusão é um condicional $\metav{A}\eif\metav{B}$, planeje usar a regra {\eif}I. Isso exige iniciar uma subprova na qual você supõe~\metav{A}. A subprova deve terminar com~\metav{B}. Em seguida, continue pensando no que deve fazer para obter $\metav{B}$ dentro dessa subprova e em como pode usar a suposição~$\metav{A}$.

Se o seu objetivo é uma conjunção, um condicional ou uma sentença negada, você deve começar trabalhando de trás para frente dessa maneira. Descreveremos em detalhe o que fazer em cada um desses casos.

\subsection*{Trabalhando de trás para frente a partir de uma conjunção}

Se queremos provar $\metav{A} \eand \metav{B}$, trabalhar de trás para frente significa escrever $\metav{A} \eand \metav{B}$ na parte inferior da nossa prova e tentar prová-la usando $\eand$I. Na parte superior, escreveremos as premissas da prova, se houver alguma. Depois, na parte inferior, escrevemos a sentença que queremos provar. Se ela for uma conjunção, vamos prová-la usando $\eand$I.
  \begin{fitchproof}
	\hypo{1}{\metav{P}_1}\PR
	\ellipsesline 
	\hypo[k]{k}{\metav{P}_k}\PR
\ellipsesline
    \have[n]{n}{\metav{A}}{} 
    \ellipsesline 
	\have[m]{m}{\metav{B}}
    \have{4}{\metav{A} \eand \metav{B}}\ai{n,m}
  \end{fitchproof}
Para $\eand$I, precisamos primeiro provar $\metav{A}$ e depois provar
$\metav{B}$. Na última linha, temos de citar as linhas nas quais
(teremos) provado $\metav{A}$ e $\metav{B}$ e usar~$\eand$I. As
partes da prova indicadas pelo $\cdots$ vertical ainda precisam ser preenchidas.
Por enquanto, marcaremos os números das linhas como $m$ e $n$. Quando a prova estiver
completa, esses marcadores poderão ser substituídos pelos números reais.

\subsection*{Trabalhando de trás para frente a partir de um condicional}

Se o nosso objetivo é provar um condicional $\metav{A} \eif \metav{B}$, teremos de usar $\eif$I. Isso exige uma subprova que comece com $\metav{A}$ e termine com~$\metav{B}$. Estruturaremos a nossa prova da seguinte forma:
\begin{fitchproof}
\open
\hypo[n]{2}{\metav{A}}
\ellipsesline 
\have[m]{3}{\metav{B}}
\close
\have{4}{\metav{A} \eif \metav{B}}\ci{2-3}
\end{fitchproof} 
Novamente, deixaremos marcadores nos espaços destinados aos números das linhas. Registraremos a última inferência como $\eif$I, citando a subprova.

\subsection*{Trabalhando de trás para frente a partir de uma sentença negada}

Se queremos provar $\enot \metav{A}$, teremos de usar $\enot$I.
\begin{fitchproof}
\open
\hypo[n]{2}{\metav{A}}
\ellipsesline 
\have[m]{3}{\ered}
\close
\have{4}{\enot \metav{A}}\ni{2-3}
\end{fitchproof} 
Para $\enot$I, temos de iniciar uma subprova com a suposição $\metav{A}$; a última linha da subprova tem de ser~$\ered$. Citaremos a subprova e usaremos~$\enot$I como regra.

Ao trabalhar de trás para frente, continue assim enquanto puder. Portanto, se você estiver trabalhando de trás para frente para provar $\metav{A} \eif \metav{B}$ e tiver estabelecido uma subprova na qual quer provar $\metav{B}$, olhe agora para~$\metav{B}$. Se, por exemplo, ele for uma conjunção, trabalhe de trás para frente a partir dele e escreva os dois conjuntivos dentro da sua subprova, e assim por diante.

\subsection*{Trabalhando de trás para frente a partir de uma disjunção}

É claro que também se pode trabalhar de trás para frente a partir de uma disjunção $\metav{A} \eor \metav{B}$, se esse for o seu objetivo.
A regra $\eor$I exige que você tenha um dos disjuntivos para inferir $\metav{A} \eor \metav{B}$.
Assim, para trabalhar de trás para frente, escolha um disjuntivo, infira $\metav{A} \eor \metav{B}$ a partir dele e, em seguida, continue procurando uma prova do disjuntivo escolhido:
\begin{fitchproof}
	\ellipsesline
	\have[n]{2}{\metav{A}} 
	\have{3}{\metav{A} \eor \metav{B}}\oi{2}
\end{fitchproof}
No entanto, talvez você não consiga provar o disjuntivo escolhido. Nesse
caso, terá de voltar atrás. Quando não conseguir preencher a parte indicada
pelo $\cdots$ vertical,
apague tudo e tente com o outro disjuntivo:
\begin{fitchproof}
	\ellipsesline 
	\have[n]{2}{\metav{B}} 
	\have{3}{\metav{A} \eor \metav{B}}\oi{2}
\end{fitchproof}
Obviamente, apagar tudo e começar de novo é frustrante, portanto você deve evitar isso. Se o seu objetivo é uma disjunção, então, você \emph{não deve começar} trabalhando de trás para frente: tente trabalhar para frente primeiro e aplique a estratégia de $\eor$I somente quando trabalhar para frente (e trabalhar de trás para frente usando $\eand$I, $\eif$I e $\enot$I) já não funcionar.

\section{Trabalhando para frente a partir do que temos}

Sua prova pode ter premissas. E, se você trabalhou de trás para frente para provar um condicional ou uma sentença negada, terá estabelecido subprovas com uma suposição e estará tentando provar uma sentença final dentro da subprova. Essas premissas e suposições são sentenças a partir das quais você pode trabalhar para frente a fim de preencher os passos ausentes da sua prova. Isso significa aplicar as regras de eliminação dos operadores principais dessas sentenças. A forma das regras dirá o que você terá de fazer.

\subsection*{Trabalhando para frente a partir de uma conjunção}

Para trabalhar para frente a partir de uma sentença da forma $\metav{A} \eand \metav{B}$, usamos $\eand$E. Essa regra nos permite fazer duas coisas: inferir $\metav{A}$ e inferir $\metav{B}$. Assim, em uma prova na qual temos $\metav{A} \eand \metav{B}$, podemos trabalhar para frente escrevendo $\metav{A}$ e/ou $\metav{B}$ imediatamente abaixo da conjunção:
\begin{fitchproof}
  \have[n]{1}{\metav{A} \eand \metav{B}}
  \have{2}{\metav{A}}\ae{1}
  \have{3}{\metav{B}}\ae{1}
\end{fitchproof}
Normalmente ficará claro, na situação particular em que você se encontra, qual de \metav{A} ou \metav{B} será necessário. No entanto, não faz mal escrever os dois.

\subsection*{Trabalhando para frente a partir de uma disjunção}

Trabalhar para frente a partir de uma disjunção funciona de maneira um pouco diferente. Para usar uma disjunção, usamos a regra $\eor$E. Para aplicar essa regra, não basta saber quais são os disjuntivos da disjunção que queremos usar. Também temos de manter em mente o que queremos provar. Suponha que queremos provar~$\metav{C}$ e temos $\metav{A} \eor B$ com que trabalhar. (Essa $\metav{A} \eor B$ pode ser uma premissa da prova, uma suposição de uma subprova ou algo já provado.) Para poder aplicar a regra $\eor$E, teremos de estabelecer duas subprovas:
\begin{fitchproof}
	\have[n]{1}{\metav{A} \eor \metav{B}}
	\open
	\hypo{2}{\metav{A}} 
	\ellipsesline 
	\have[m]{3}{\metav{C}}
	\close 
	\open
	\hypo{4}{\metav{B}}
	\ellipsesline
	\have[k]{5}{\metav{C}}
	\close
	\have{6}{\metav{C}}\oe{1,(2)-3,(4)-5} 
\end{fitchproof} 
A primeira subprova começa com o primeiro disjuntivo, $\metav{A}$, e termina com a sentença que estamos procurando, $\metav{C}$. A segunda subprova começa com o outro disjuntivo, $\metav{B}$, e também termina com a sentença objetivo~$\metav{C}$. Cada uma dessas subprovas ainda precisa ser preenchida. Então, podemos justificar a sentença objetivo $\metav{C}$ usando $\eor$E, citando a linha com $\metav{A} \eor \metav{B}$ e as duas subprovas.

\subsection*{Trabalhando para frente a partir de um condicional}

Para usar um condicional $\metav{A} \eif \metav{B}$, você também precisa do antecedente $\metav{A}$ para aplicar~$\eif$E. Assim, para trabalhar para frente a partir de um condicional, você derivará $\metav{B}$, justificá-lo-á por $\eif$E e estabelecerá $\metav{A}$ como um novo subobjetivo.
\begin{fitchproof}
	\have[n]{1}{\metav{A} \eif \metav{B}}
	\ellipsesline 
	\have[m]{2}{\metav{A}}
	\have{3}{\metav{B}}\ce{1,2} 
\end{fitchproof}

\subsection*{Trabalhando para frente a partir de uma sentença negada}

Por fim, para usar uma sentença negada $\enot \metav{A}$, você aplicaria $\enot$E. Além de $\enot \metav{A}$, ela exige também a sentença correspondente~$\metav{A}$ sem a negação. A sentença que você obterá é sempre a mesma: $\ered$. Assim, trabalhar para frente a partir de uma sentença negada funciona especialmente bem dentro de uma subprova que você pretende usar para $\enot$I (ou IP). Você trabalha para frente a partir de $\enot \metav{A}$ se já tem $\enot \metav{A}$ e quer provar~$\ered$. Para isso, estabeleça $\metav{A}$ como um novo subobjetivo.
\begin{fitchproof}
	\have[n]{1}{\enot \metav{A}}
	\ellipsesline 
	\have[m]{2}{\metav{A}}
	\have{3}{\ered}\ne{1,2} 
\end{fitchproof}

\section{Estratégias em ação}

Suponha que queremos mostrar que o argumento $(A \eand B) \eor (A \eand C) \therefore A \eand (B \eor C)$ é válido. Começamos a prova escrevendo a premissa e a conclusão. (Em uma folha de papel, você vai querer deixar o máximo de espaço possível entre elas; portanto, escreva as premissas no alto da folha e a conclusão na parte inferior.)
\begin{fitchproof}
   \hypo{1}{(A \eand B) \eor (A \eand C)}\PR
\ellipsesline
  \have[n]{2}{A \eand (B \eor C)}
\end{fitchproof}
Agora temos duas opções: trabalhar de trás para frente a partir da conclusão ou trabalhar para frente a partir da premissa. Escolheremos a segunda estratégia: usamos a disjunção na linha~$1$ e estabelecemos as subprovas necessárias para $\eor$E. A disjunção na linha~$1$ tem dois disjuntivos, $A \eand B$ e $A \eand C$. A sentença objetivo que você quer provar é $A \eand (B \eor C)$. Portanto, neste caso, você tem de estabelecer duas subprovas: uma com a suposição $A \eand B$ e a última linha $A \eand (B \eor C)$; a outra com a suposição $A \eand C$ e a última linha $A \eand (B \eor C)$. A justificativa para a conclusão na linha~$n$ será $\eor$E, citando a disjunção na linha~$1$ e as duas subprovas. Assim, sua prova passa a ter a seguinte aparência:
\begin{fitchproof}
	\hypo{1}{(A \eand B) \eor (A \eand C)}\PR
	\open
	\hypo{2}{A \eand B}\AS
	\ellipsesline 
	\have[n]{6}{A \eand (B \eor C)}
	\close
	\open
	\hypo{7}{A \eand C}\AS
	\ellipsesline
	\have[m]{11}{A \eand (B \eor C)}
	\close
	\have{12}{A \eand (B \eor C)}\oe{1,2-6,(7)-11}
\end{fitchproof}
Agora você tem duas tarefas separadas, a saber, preencher cada uma das duas subprovas. Na primeira subprova, trabalhamos de trás para frente a partir da conclusão $A \eand (B \eor C)$. Essa é uma conjunção; portanto, dentro da primeira subprova, você terá dois subobjetivos separados: provar $A$ e provar $B \eor C$. Esses subobjetivos permitirão justificar a linha~$n$ usando~$\eand$I. Sua prova passa a ter a seguinte aparência:
\begin{fitchproof}
	\hypo{1}{(A \eand B) \eor (A \eand C)}\PR
	\open
	\hypo{2}{A \eand B}\AS
	\ellipsesline
	\have[i]{4}{A}
	\ellipsesline
	\have[n][-1]{5}{B \eor C}
	\have[n]{6}{A \eand (B \eor C)}\ai{4,5}
	\close
	\open
	\hypo{7}{A \eand C}\AS
	\ellipsesline
	\have[m]{11}{A \eand (B \eor C)}
	\close
	\have{12}{A \eand (B \eor C)}\oe{1,2-6,(7)-11}
\end{fitchproof}
Vemos imediatamente que podemos obter a linha $i$ a partir da linha~$2$ por $\eand$E. Portanto, a linha $i$ é, na verdade, a linha~$3$ e pode ser justificada com $\eand$E a partir da linha~$2$. O outro subobjetivo, $B \eor C$, é uma disjunção. Aplicaremos a estratégia de trabalhar de trás para frente a partir de uma disjunção até a linha $n-1$. Temos a opção de escolher $B$ ou~$C$ como subobjetivo. Escolher $C$ não funcionaria e acabaríamos tendo de voltar atrás. E você já pode ver que, se escolher $B$ como subobjetivo, poderá obtê-lo trabalhando novamente para frente a partir da conjunção $A \eand B$ na linha~$2$. Assim, podemos completar a primeira subprova da seguinte maneira:
\begin{fitchproof}
	\hypo{1}{(A \eand B) \eor (A \eand C)}\PR
	\open
	\hypo{2}{A \eand B}\AS
	\have{3}{A}\ae{2}
	\have{4}{B}\ae{2}
	\have{5}{B \eor C}\oi{4}
	\have{6}{A \eand (B \eor C)}\ai{3,5}
	\close
	\open
	\hypo{7}{A \eand C}\AS
	\ellipsesline
	\have[m]{11}{A \eand (B \eor C)}
	\close
	\have{12}{A \eand (B \eor C)}\oe{1,2-6,7-11}
\end{fitchproof}
Assim como a linha~$3$, obtemos a linha~$4$ a partir da linha~$2$ por $\eand$E. A linha~$5$ é justificada por $\eor$I a partir da linha~$4$, pois ali estávamos trabalhando de trás para frente a partir de uma disjunção.

Isso conclui a primeira subprova. A segunda subprova é quase exatamente igual. Deixaremos isso como exercício.

Lembre-se de que, quando começamos, tínhamos a opção de trabalhar para frente a partir da premissa ou de trabalhar de trás para frente a partir da conclusão, e escolhemos a primeira opção. A segunda opção também conduz a uma prova, mas ela terá uma aparência diferente. Os primeiros passos seriam trabalhar de trás para frente a partir da conclusão, estabelecendo dois subobjetivos, $A$ e $B \eor C$, e então trabalhar para frente a partir da premissa para prová-los, por exemplo:
\begin{fitchproof}
	\hypo{1}{(A \eand B) \eor (A \eand C)}\PR
	\open
	\hypo{2}{A \eand B}\AS
	\ellipsesline
	\have[k]{3}{A}
	\close
	\open
	\hypo{4}{A \eand C}\AS
	\ellipsesline
	\have[n][-1]{5}{A}
	\close
	\have{6}{A}\oe{1,2-3,(4)-(5)}
	\open
	\hypo{7}{A \eand B}\AS
	\ellipsesline
	\have[l]{8}{B \eor C}
	\close
	\open
	\hypo{9}{A \eand C}\AS
	\ellipsesline
	\have[m][-1]{10}{B \eor C}
	\close
	\have{11}{B \eor C}\oe{1,(7)-8,(9)-(10)}	
	\have{12}{A \eand (B \eor C)}\ai{6,11}
\end{fitchproof}
Deixaremos para você o preenchimento das partes ausentes indicadas por~$\vdots$.

Vejamos outro exemplo para ilustrar como aplicar as estratégias
para lidar com condicionais e negação. A sentença $(A \eif B) \eif
(\enot B \eif \enot A)$ é uma tautologia. Vamos ver se conseguimos encontrar uma
prova dela, sem premissas, usando as estratégias. Primeiro escrevemos
a sentença na parte inferior de uma folha de papel. Como trabalhar para frente
não é uma opção (não há nada a partir do qual trabalhar para frente), trabalhamos
de trás para frente e estabelecemos uma subprova para provar a sentença desejada, $(A
\eif B) \eif (\enot B \eif \enot A)$, usando $\eif$I. Sua suposição
deve ser o antecedente do condicional que queremos provar, isto é, $A
\eif B$, e sua última linha deve ser o consequente, $\enot B \eif \enot A$.
\begin{fitchproof}
\open
\hypo{1}{A \eif B}\AS
\ellipsesline
\have[n]{7}{\enot B \eif \enot A}
\close
\have{8}{(A \eif B) \eif (\enot B \eif \enot A)}\ci{1-7}
\end{fitchproof}
O novo objetivo, $\enot B \eif \enot A$, também é um condicional; portanto, trabalhando de trás para frente, estabelecemos outra subprova:
\begin{fitchproof}
	\open
	\hypo{1}{A \eif B}\AS
	\open
	\hypo{2}{\enot B}\AS
	\ellipsesline
	\have[n][-1]{6}{\enot A}
	\close
	\have{7}{\enot B \eif \enot A}\ci{2-(6)}
	\close
	\have{8}{(A \eif B) \eif (\enot B \eif \enot A)}\ci{1-7}
\end{fitchproof}
A partir de $\enot A$, trabalhamos novamente de trás para frente. Para isso, olhe para a regra $\enot$I. Ela exige uma subprova com~$A$ como suposição e $\ered$ como sua última linha. Assim, a prova agora é:
\begin{fitchproof}
	\open
	\hypo{1}{A \eif B}\AS
	\open
	\hypo{2}{\enot B}\AS
	\open\hypo{3}{A}
	\ellipsesline
	\have[n][-2]{5}{\ered}
	\close
	\have{6}{\enot A}\ni{3-(5)}
	\close
	\have{7}{\enot B \eif \enot A}\ci{2-(6)}
	\close
	\have{8}{(A \eif B) \eif (\enot B \eif \enot A)}\ci{1-7}
\end{fitchproof}
Agora o nosso objetivo é provar~$\ered$. Como dissemos acima, ao discutir
como trabalhar para frente a partir de uma sentença negada, a regra $\enot$E permite
provar~$\ered$, que é o nosso objetivo na subprova mais interna. Assim,
procuramos uma sentença negada a partir da qual possamos trabalhar para frente:
seria $\enot B$ na linha~$2$. Isso significa que temos de derivar $B$
dentro da subprova, pois $\enot$E exige não apenas $\enot B$ (que já
temos), mas também~$B$. E, por sua vez, obtemos $B$ trabalhando para
frente a partir de $A \eif B$, pois $\eif$E nos permitirá justificar o
consequente desse condicional, $B$, por $\eif$E. A regra $\eif$E também
exige o antecedente~$A$ do condicional, mas ele também já está
disponível (na linha~$3$). Portanto, terminamos com:
\begin{fitchproof}
	\open
	\hypo{1}{A \eif B}\AS
	\open
	\hypo{2}{\enot B}\AS
	\open\hypo{3}{A}
	\have{4}{B}\ce{1,3}
	\have{5}{\ered}\ne{2,4}
	\close
	\have{6}{\enot A}\ni{3-5}
	\close
	\have{7}{\enot B \eif \enot A}\ci{2-6}
	\close
	\have{8}{(A \eif B) \eif (\enot B \eif \enot A)}\ci{1-7}
\end{fitchproof}

\section{Trabalhando para frente a partir de $\ered$}\label{sec:backred}

Ao aplicar as estratégias, às vezes você se encontrará em uma situação em que pode justificar~$\ered$. Usando a regra da explosão, isso permitiria justificar \emph{qualquer coisa}. Assim, $\ered$ funciona como um curinga nas provas. Por exemplo, suponha que você queira dar uma prova do argumento $A \eor B, \enot A \therefore B$. Você estrutura a prova escrevendo as premissas $A \eor B$ e $\enot A$ no alto, nas linhas $1$ e $2$, e a conclusão~$B$ na parte inferior da página. $B$ não tem conectivo principal, então não se pode trabalhar de trás para frente a partir dele. Em vez disso, é preciso trabalhar para frente a partir de $A \eor B$: isso exige duas subprovas, assim:
\begin{fitchproof}
	\hypo{1}{A \eor B}\PR
	\hypo{7}{\enot A}\PR
	\open
	\hypo{2}{A} \AS
	\ellipsesline 
	\have[m]{3}{B}
	\close 
	\open
	\hypo{4}{B}\AS
	\ellipsesline
	\have[k]{5}{B}
	\close
	\have{6}{B}\oe{1,2-3,(4)-5} 
\end{fitchproof} 
Observe que você tem $\enot A$ na linha~$2$ e $A$ como suposição da sua primeira subprova. Isso lhe dá $\ered$ usando $\enot$E e, a partir de $\ered$, você obtém a conclusão~$B$ da primeira subprova usando~X. Lembre-se de que é possível repetir uma sentença que você já tem usando a regra de reiteração~R. Portanto, a nossa prova seria:
\begin{fitchproof}
	\hypo{1}{A \eor B}\PR
	\hypo{7}{\enot A}\PR
	\open
	\hypo{2}{A} \AS
	\have{8}{\ered}\ne{7,2} 
	\have{3}{B}\re{8}
	\close 
	\open
	\hypo{4}{B}\AS
	\have{5}{B}\by{R}{4}
	\close
	\have{6}{B}\oe{1,2-3,4-5} 
\end{fitchproof} 

\section{Procedendo indiretamente}

Em muitos casos, as estratégias de trabalhar para frente e de trás para frente acabam funcionando. Mas há casos em que elas não funcionam. Se você não conseguir encontrar uma maneira de mostrar $\metav{A}$ diretamente usando essas estratégias, use IP. Para isso, estabeleça uma subprova na qual você supõe $\enot\metav{A}$ e procure uma prova de $\ered$ dentro dessa subprova.

\begin{fitchproof}
\open
\hypo[n]{2}{\enot \metav{A}}\AS
\ellipsesline 
\have[m]{3}{\ered}
\close
\have{4}{\metav{A}}\ip{2-3}
\end{fitchproof}
Aqui, temos de iniciar uma subprova com a suposição $\enot \metav{A}$;
a última linha da subprova tem de ser~$\ered$. Citaremos a subprova e usaremos~IP como regra. Na subprova, agora temos uma suposição adicional (na linha $n$) com a qual trabalhar.

Suponha que usamos a estratégia da prova indireta ou que estamos em alguma outra situação na qual procuramos uma prova de $\ered$. Qual é um bom candidato? Naturalmente, o candidato óbvio seria usar uma sentença negada, pois, como vimos acima, $\enot$E sempre produz~$\ered$. Se você estrutura uma prova como a anterior, tentando provar \metav{A} usando~IP, terá $\enot \metav{A}$ como suposição da sua subprova; portanto, trabalhando para frente a partir dela para justificar $\ered$ dentro da subprova, estabelecerá em seguida \metav{A} como um objetivo dentro da subprova. Se estiver usando essa estratégia IP, você se encontrará na seguinte situação:
\begin{fitchproof}
\open
\hypo[n]{2}{\enot \metav{A}}\AS
\ellipsesline
\have[m][-1]{3}{\metav{A}}
\have{4}{\ered}\ne{2,3}
\close
\have{5}{\metav{A}}\ip{2-4}
\end{fitchproof} 
Isso parece estranho: queríamos provar $\metav{A}$ e as estratégias falharam; por isso, usamos IP como último recurso. E agora nos encontramos na mesma situação: estamos novamente procurando uma prova de~$\metav{A}$. Mas observe que agora estamos \emph{dentro} de uma subprova e, nessa subprova, temos uma suposição adicional ($\enot \metav{A}$) com a qual trabalhar, algo que não tínhamos antes. Vejamos um exemplo.

\section{Prova indireta do terceiro excluído}\label{s:proofLEM}

A sentença $A \eor \enot A$ é uma tautologia chamada \emph{lei do
terceiro excluído} (LTE), pois expressa a ideia de que \metav{A} pode ser verdadeira ou $\enot \metav{A}$ pode ser verdadeira, mas não há uma via intermediária em que nenhuma das duas seja verdadeira.\footnote{Às vezes, você pode encontrar lógicos ou filósofos falando em ``tertium non datur''. Esse é o mesmo princípio do terceiro excluído; significa ``não há uma terceira via''. Lógicos que têm reservas quanto à prova indireta também têm reservas quanto à LTE.} Ela deve ter uma prova mesmo sem premissas. Mas trabalhar de trás para frente não nos ajuda: para obter $A \eor \enot A$ usando $\eor$I, teríamos de provar $A$ ou $\enot A$---novamente, sem premissas. Nenhuma das duas é uma tautologia, portanto não conseguiremos provar nenhuma delas. Trabalhar para frente também não funciona, pois não há nada a partir do qual trabalhar para frente. Assim, a única opção é a prova indireta.
\begin{fitchproof}
	\open
	\hypo{2}{\enot (A \eor \enot A)}\AS
	\ellipsesline
	\have[m]{8}{\ered}
	\close
	\have{9}{A \eor \enot A}\ip{2-8}
\end{fitchproof}
Agora temos algo a partir do qual trabalhar para frente: a suposição $\enot(A \eor \enot A)$. Para usá-la, justificamos $\ered$ por $\enot$E, citando a suposição na linha~$1$ e também a sentença correspondente sem negação, $A \eor \enot A$, que ainda precisa ser provada.
\begin{fitchproof}
	\open
	\hypo{2}{\enot (A \eor \enot A)}\AS
	\ellipsesline
	\have[m][-1]{7}{A \eor \enot A}
	\have{8}{\ered}\ne{2,7}
	\close
	\have{9}{A \eor \enot A}\ip{2-8}
\end{fitchproof}
No início, trabalhar de trás para frente para provar $A \eor\enot A$ por $\eor$I não funcionou. Mas agora estamos em uma situação diferente: queremos provar $A \eor\enot A$ dentro de uma subprova. Em geral, ao lidar com novos objetivos, devemos voltar e começar com as estratégias básicas. Neste caso, devemos primeiro tentar trabalhar de trás para frente a partir da disjunção $A \eor \enot A$, isto é, temos de escolher um disjuntivo e tentar prová-lo. Vamos escolher~$\enot A$. Isso nos permitiria justificar $A \eor \enot A$ na linha~$m - 1$ usando $\eor$I. Então, trabalhando de trás para frente a partir de $\enot A$, iniciamos outra subprova para justificar $\enot A$ usando $\enot$I. Essa subprova deve ter $A$ como suposição e~$\ered$ como sua última linha.
\begin{fitchproof}
	\open
	\hypo{2}{\enot (A \eor \enot A)}\AS
	\open
	\hypo{3}{A}\AS
	\ellipsesline
	\have[m][-3]{5}{\ered}
	\close
	\have{6}{\enot A}\ni{3-(5)}
	\have{7}{A \eor \enot A}\oi{6}
	\have{8}{\ered}\ne{2,7}
	\close
	\have{9}{A \eor \enot A}\ip{2-8}
\end{fitchproof}
Dentro dessa nova subprova, novamente precisamos justificar $\ered$. A melhor maneira de fazer isso é trabalhar para frente a partir de uma sentença negada; $\enot(A \eor \enot A)$ na linha~$1$ é a única sentença negada que podemos usar. A sentença correspondente sem negação, $A \eor \enot A$, porém, segue diretamente de $A$ (que temos na linha~$2$) por $\eor$I. A nossa prova completa é:
\begin{fitchproof}
	\open
	\hypo{2}{\enot (A \eor \enot A)}\AS
	\open
	\hypo{3}{A}\AS
	\have{4}{A \eor \enot A}\oi{3}
	\have{5}{\ered}\ne{2,4}
	\close
	\have{6}{\enot A}\ni{3-5}
	\have{7}{A \eor \enot A}\oi{6}
	\have{8}{\ered}\ne{2,7}
	\close
	\have{9}{A \eor \enot A}\ip{2-8}
\end{fitchproof}

\practiceproblems

\problempart
Use as estratégias para encontrar provas para cada um dos argumentos a seguir:
\begin{compactlist}
\item $A \eif B, A \eif C \therefore A \eif (B \eand C)$
\item $(A \eand B) \eif C \therefore A \eif (B \eif C)$
\item $A \eif (B \eif C) \therefore (A \eif B) \eif (A \eif C)$
\item $A \eor (B \eand C) \therefore (A \eor B) \eand (A \eor C)$
\item $(A \eand B) \eor (A \eand C) \therefore A \eand (B \eor C)$
\item $A \eor B, A \eif C, B \eif D \therefore C \eor D$
\item $\enot A \lor \enot B \therefore \enot(A \eand B)$
\item $A \eand \enot B \therefore \enot(A \eif B)$
\end{compactlist}

\problempart
Formule estratégias para trabalhar de trás para frente e para frente a partir de $\metav{A} \eiff \metav{B}$.

\problempart
Use as estratégias para encontrar provas para cada uma das sentenças a seguir:
\begin{compactlist}
\item $\enot A \eif (A \eif \ered)$
\item $\enot(A \eand \enot A)$
\item $[(A \eif C) \eand (B \eif C)] \eif [(A \lor B) \eif C]$
\item $\enot(A \eif B) \eif (A \eand \enot B)$
\item $(\enot A \eor B) \eif (A \eif B)$
\end{compactlist}
Como estas devem ser provas de sentenças sem premissas, você começará com a respectiva sentença na \emph{parte inferior} da prova, que não terá premissas.

\problempart
Use as estratégias para encontrar provas para cada um dos argumentos e sentenças a seguir:
\begin{compactlist}
\item $\enot\enot A \eif A$
\item $\enot A \eif \enot B \therefore B \eif A$
\item $A \eif B \therefore \enot A \eor B$
\item $\enot(A \eand B) \eif (\enot A \eor \enot B)$
\item $A \eif (B \eor C) \therefore (A \eif B) \eor (A \eif C)$
\item $(A \eif B) \lor (B \eif A)$
\item $((A \eif B) \eif A) \eif A$
\end{compactlist}
Todos exigirão a estratégia IP. Os três últimos são especialmente difíceis!

\chapter{Regras adicionais para a LVF}\label{s:Further}
Em \cref{s:BasicTFL}, introduzimos as regras básicas do nosso sistema de provas para a LVF. Nesta seção, acrescentaremos algumas regras ao nosso sistema. O sistema de provas estendido é um pouco mais fácil de usar. (Contudo, em \cref{s:Derived}, veremos que, estritamente falando, elas não são \emph{necessárias}.)

% \section{Reiteration}
% The first additional rule is \emph{reiteration} (R). This just allows us to repeat ourselves:
% \factoidbox{\begin{fitchproof}
% 	\have[m]{a}{\metav{A}}
% 	\have[\ ]{b}{\metav{A}} \by{R}{a}
% \end{fitchproof}}
% Such a rule is obviously legitimate; we could have used it in our proof in \cref{sec:backred}:
% \begin{fitchproof}
% 	\hypo{1}{A \eor B}
% 	\hypo{7}{\enot A}
% 	\open
% 	\hypo{2}{A} 
% 	\have{8}{\ered}\ne{7,2} 
% 	\have{3}{B}\re{8}
% 	\close 
% 	\open
% 	\hypo{4}{B}
% 	\have{5}{B}\by{R}{4}
% 	\close
% 	\have{6}{B}\oe{1,2-3,4-5} 
% \end{fitchproof}
% This is a fairly typical use of the R rule.

\section{Silogismo disjuntivo}
Eis uma forma de argumento muito natural.
	\begin{earg}
\item Elizabeth está em Massachusetts ou em DC.
\item Ela não está em DC.
\item[\texttherefore] Ela está em Massachusetts.
	\end{earg}
Isso é chamado de \emph{silogismo disjuntivo}. Nós o adicionamos ao nosso sistema de provas da seguinte forma:
\factoidbox{\begin{fitchproof}
	\have[m]{ab}{\metav{A} \eor \metav{B}}
	\have[n]{nb}{\enot \metav{A}}
	\have[\ ]{con}{\metav{B}}\by{DS}{ab, nb}
\end{fitchproof}}
and
\factoidbox{\begin{fitchproof}
	\have[m]{ab}{\metav{A} \eor \metav{B}}
	\have[n]{nb}{\enot \metav{B}}
	\have[\ ]{con}{\metav{A}}\by{DS}{ab, nb}
\end{fitchproof}}
Como de costume, a disjunção e a negação de um dos disjuntivos podem ocorrer em qualquer ordem e não precisam ser adjacentes. No entanto, sempre citamos primeiro a disjunção.

\section{Modus tollens}
Outro padrão útil de inferência aparece no argumento a seguir:
	\begin{earg}
		\item Se Hilary venceu a eleição, então ela está na Casa Branca.
		\item Ela não está na Casa Branca.
		\item[\texttherefore] Ela não venceu a eleição.
	\end{earg}
Esse padrão de inferência é chamado de \emph{modus tollens}. A regra correspondente é:
\factoidbox{\begin{fitchproof}
	\have[m]{ab}{\metav{A}\eif\metav{B}}
	\have[n]{a}{\enot\metav{B}}
	\have[\ ]{b}{\enot\metav{A}}\mt{ab,a}
\end{fitchproof}}
Como de costume, as premissas podem ocorrer em qualquer ordem, mas sempre citamos primeiro o condicional.

\section{Eliminação da dupla negação}
Outra regra útil é a \emph{eliminação da dupla negação}. Ela faz exatamente o que seu nome indica:
\factoidbox{\begin{fitchproof}
		\have[m]{dna}{\enot \enot \metav{A}}
		\have[ ]{a}{\metav{A}}\dne{dna}
	\end{fitchproof}}
A justificativa para isso é que, na linguagem natural, duplas negações tendem a se cancelar.

Dito isso, você deve estar ciente de que o contexto e a ênfase podem impedir que isso aconteça. Considere: \textquotedblleft{}Jane não está \emph{não} feliz\textquotedblright{}. Pode-se argumentar que não é possível inferir \textquotedblleft{}Jane está feliz\textquotedblright{}, pois a primeira sentença deve ser entendida como equivalente a \textquotedblleft{}Jane não está \emph{infeliz}\textquotedblright{}. Isso é compatível com \textquotedblleft{}Jane está em um estado de profunda indiferença\textquotedblright{}. Como de costume, passar para a LVF nos obriga a sacrificar certas nuances das expressões em linguagem natural.

\section{Terceiro excluído}

Suponha que possamos mostrar que, se estiver ensolarado lá fora, Bill terá levado um guarda-chuva (por medo de se queimar). Suponha também que possamos mostrar que, se não estiver ensolarado lá fora, Bill terá levado um guarda-chuva (por medo da chuva). Bem, não há uma terceira possibilidade para o tempo. Portanto, \emph{qualquer que seja o tempo}, Bill terá levado um guarda-chuva.

Essa linha de raciocínio motiva a regra a seguir:
\factoidbox{\begin{fitchproof}
		\open
			\hypo[i]{a}{\metav{A}}\AS
			\have[j]{c1}{\metav{B}}
		\close
		\open
			\hypo[k]{b}{\enot\metav{A}}\AS
			\have[l]{c2}{\metav{B}}
		\close
		\have[\ ]{ab}{\metav{B}}\tnd{a-c1,b-c2}
	\end{fitchproof}}
Às vezes, a regra é chamada de \emph{regra do terceiro excluído}. Você
pode entendê-la como $\eor$E aplicado a $\metav{A} \eor \enot \metav{A}$.
Pode haver quantas linhas você quiser entre $i$ e $j$ e quantas linhas
você quiser entre $k$ e $l$. Além disso, as subprovas podem aparecer em
qualquer ordem, e a segunda subprova não precisa vir imediatamente após a primeira.

Para ver a regra em ação, considere:
	$$P \therefore (P \eand D) \eor (P \eand \enot D)$$
Eis uma prova correspondente ao argumento:
	\begin{fitchproof}
		\hypo{a}{P}\PR
		\open
			\hypo{b}{D}\AS
			\have{ab}{P \eand D}\ai{a, b}
			\have{abo}{(P \eand D) \eor (P \eand \enot D)}\oi{ab}
		\close
		\open
			\hypo{nb}{\enot D}\AS
			\have{anb}{P \eand \enot D}\ai{a, nb}
			\have{anbo}{(P \eand D) \eor (P \eand \enot D)}\oi{anb}
		\close
		\have{con}{(P \eand D) \eor (P \eand \enot D)}\tnd{b-abo, nb-anbo}
	\end{fitchproof}
Eis outro exemplo:
\begin{fitchproof}
	\hypo{ana}{A \eif \enot A}\PR
	\open
		\hypo{a}{A}\AS
		\have{na}{\enot A}\ce{ana, a}
	\close
	\open
		\hypo{na1}{\enot A}\AS
		\have{na2}{\enot A}\by{R}{na1}
	\close
	\have{na3}{\enot A}\tnd{a-na, na1-na2}
\end{fitchproof}


\section{Regras de De Morgan}
As nossas últimas regras adicionais são chamadas de leis de De~Morgan (em homenagem a Augustus De~Morgan). A forma das regras deve ser familiar a partir das tabelas-verdade.

A primeira regra de De Morgan é:
\factoidbox{\begin{fitchproof}
	\have[m]{ab}{\enot (\metav{A} \eand \metav{B})}
	\have[\ ]{dm}{\enot \metav{A} \eor \enot \metav{B}}\dem{ab}
\end{fitchproof}}
A segunda regra de De Morgan é o inverso da primeira:
\factoidbox{\begin{fitchproof}
	\have[m]{ab}{\enot \metav{A} \eor \enot \metav{B}}
	\have[\ ]{dm}{\enot (\metav{A} \eand \metav{B})}\dem{ab}
\end{fitchproof}}
A terceira regra de De Morgan é a \emph{dual} da primeira:
\factoidbox{\begin{fitchproof}
	\have[m]{ab}{\enot (\metav{A} \eor \metav{B})}
	\have[\ ]{dm}{\enot \metav{A} \eand \enot \metav{B}}\dem{ab}
\end{fitchproof}}
E a quarta é o inverso da terceira:
\factoidbox{\begin{fitchproof}
	\have[m]{ab}{\enot \metav{A} \eand \enot \metav{B}}
	\have[\ ]{dm}{\enot (\metav{A} \eor \metav{B})}\dem{ab}
\end{fitchproof}}
\emph{Essas são todas as regras adicionais do nosso sistema de provas para a LVF.}

\practiceproblems
\solutions
\problempart
\label{pr.justifyTFLproof}
As provas a seguir estão sem suas citações (regras e números de linhas). Acrescente-as onde forem necessárias:
\begin{compactlist}
\item\begin{fitchproof}
\hypo{1}{W \eif \enot B}
\hypo{2}{A \eand W}
\hypo{2b}{B \eor (J \eand K)}
\have{3}{W}{}
\have{4}{\enot B} {}
\have{5}{J \eand K} {}
\have{6}{K}{}
\end{fitchproof}
\item\begin{fitchproof}
\hypo{1}{L \eiff \enot O}
\hypo{2}{L \eor \enot O}
\open
	\hypo{a1}{\enot L}
	\have{a2}{\enot O}{}
	\have{a3}{L}{}
	\have{a4}{\ered}{}
\close
\have{3b}{\enot\enot L}{}
\have{3}{L}{}
\end{fitchproof}
\item\begin{fitchproof}
\hypo{1}{Z \eif (C \eand \enot N)}
\hypo{2}{\enot Z \eif (N \eand \enot C)}
\open
	\hypo{a1}{\enot(N \eor  C)}
	\have{a2}{\enot N \eand \enot C} {}
	\have{a6}{\enot N}{}
	\have{b4}{\enot C}{}
		\open
		\hypo{b1}{Z}
		\have{b2}{C \eand \enot N}{}
		\have{b3}{C}{}
		\have{red}{\ered}{}
	\close
	\have{a3}{\enot Z}{}
	\have{a4}{N \eand \enot C}{}
	\have{a5}{N}{}
	\have{a7}{\ered}{}
\close
\have{3b}{\enot\enot(N \eor C)}{}
\have{3}{N \eor C}{}
\end{fitchproof}
\end{compactlist}

\problempart 
Forneça uma prova para cada um destes argumentos:
\begin{compactlist}
\item $E\eor F$, $F\eor G$, $\enot F \therefore E \eand G$
\item $M\eor(N\eif M) \therefore \enot M \eif \enot N$
\item $(M \eor N) \eand (O \eor P)$, $N \eif P$, $\enot P \therefore M\eand O$
\item $(X\eand Y)\eor(X\eand Z)$, $\enot(X\eand D)$, $D\eor M \therefore M$
\end{compactlist}



\chapter{Conceitos prova-teóricos}\label{s:ProofTheoreticConcepts}

Neste capítulo, introduziremos algum vocabulário novo. A expressão a seguir:
$$\metav{A}_1, \metav{A}_2, \ldots, \metav{A}_n \proves \metav{C}$$
significa que existe alguma prova que termina com $\metav{C}$ cujas suposições não descarregadas estão entre $\metav{A}_1, \metav{A}_2, \ldots, \metav{A}_n$. Quando queremos dizer que \emph{não} é o caso de existir uma prova que termine com $\metav{C}$ a partir de $\metav{A}_1$, $\metav{A}_2$, \dots,~$\metav{A}_n$, escrevemos:  $$\metav{A}_1, \metav{A}_2, \ldots, \metav{A}_n \nproves \metav{C}$$

O símbolo \textquotedblleft{}$\proves$\textquotedblright{} é chamado de \emph{turnstile simples}. Queremos enfatizar que ele não é o símbolo de {turnstile duplo} (\textquotedblleft{}$\entails$\textquotedblright{}) que introduzimos em \cref{s:SemanticConcepts} para simbolizar a consequência lógica. O turnstile simples, \textquotedblleft{}$\proves$\textquotedblright{}, diz respeito à existência de provas; o turnstile duplo, \textquotedblleft{}$\entails$\textquotedblright{}, diz respeito à existência de valorações (ou interpretações, quando usado para a LPO). \emph{São noções muito diferentes.}

Munidos do símbolo \textquotedblleft{}$\proves$\textquotedblright{}, podemos introduzir mais terminologia. Para dizer que existe uma prova de $\metav{A}$ sem suposições não descarregadas, escrevemos: ${} \proves \metav{A}$. Nesse caso, dizemos que $\metav{A}$ é um \define{teorema}.
	\factoidbox{\label{def:syntactic_tautology_in_sl}
		$\metav{A}$ é um \define{teorema} \ifeff{} $\proves \metav{A}$
	}
\newglossaryentry{theorem}
{
name=teorema,
description={Uma senten\c{c}a que pode ser provada sem quaisquer premissas}
}

Para ilustrar isso, suponha que queremos mostrar que \textquotedblleft{}$\enot (A \eand \enot A)$\textquotedblright{} é um teorema. Portanto, precisamos de uma prova de \textquotedblleft{}$\enot(A \eand \enot A)$\textquotedblright{} que não tenha \emph{nenhuma} suposição não descarregada. Contudo, como queremos provar uma sentença cujo operador lógico principal é uma negação, vamos começar com uma \emph{subprova} dentro da qual supomos \textquotedblleft{}$A \eand \enot A$\textquotedblright{} e mostrar que essa suposição conduz a uma contradição. Em suma, a prova é a seguinte:
	\begin{fitchproof}
		\open
			\hypo{con}{A \eand \enot A}\AS
			\have{a}{A}\ae{con}
			\have{na}{\enot A}\ae{con}
			\have{red}{\ered}\ne{na, a}
		\close
		\have{lnc}{\enot (A \eand \enot A)}\ni{con-red}
	\end{fitchproof}
Portanto, provamos \textquotedblleft{}$\enot (A \eand \enot A)$\textquotedblright{} sem
suposições (não descarregadas). Esse teorema particular é um caso
do que às vezes se chama \emph{lei da não contradição}.

Para mostrar que algo é um teorema, basta encontrar uma prova adequada. Normalmente, é muito mais difícil mostrar que algo \emph{não} é um teorema. Para isso, seria preciso demonstrar não apenas que determinadas estratégias de prova falham, mas que \emph{nenhuma} prova é possível. Mesmo que você não consiga provar uma sentença de mil maneiras diferentes, talvez a prova seja simplesmente longa e complexa demais para que consiga compreendê-la. Talvez você simplesmente não tenha tentado o suficiente.

Eis outro termo novo:
	\factoidbox{
		Duas sentenças \metav{A} e \metav{B} são \define{interderiv\'aveis} \ifeff{} cada uma pode ser provada a partir da outra; isto é, tanto $\metav{A}\proves\metav{B}$ quanto $\metav{B}\proves\metav{A}$.
	}
        
\newglossaryentry{interderivable}
{
  name=interderivabilidade,
  text = interderiv\'avel,
description={Uma propriedade de pares de senten\c{c}as se, e somente se, existe uma deriva\c{c}\~{a}o que leva de cada uma \`a outra}
}


Assim como no caso de mostrar que uma sentença é um teorema, é relativamente fácil mostrar que duas sentenças são interderiváveis: basta fornecer um par de provas. Mostrar que sentenças \emph{não} são interderiváveis seria muito mais difícil: é tão difícil quanto mostrar que uma sentença não é um teorema.

Eis um terceiro termo relacionado:
	\factoidbox{
	As sentenças $\metav{A}_1,\metav{A}_2,\ldots, \metav{A}_n$ são
	\define{conjuntamente inconsistentes} \ifeff{} a contradição~$\ered$ pode ser provada a partir
	delas, isto é, $\metav{A}_1,\metav{A}_2,\ldots, \metav{A}_n \proves
	\ered$. Se elas não são \define{inconsistentes}, nós as chamamos de
	\define{conjuntamente consistentes}.
	}
        
\newglossaryentry{inconsistent}
{    name={inconsist\^encia},
  description={Senten\c{c}as s\~{a}o inconsistentes \ifeff{} a contradi\c{c}\~{a}o~$\ered$ pode ser provada a partir delas},
    text={inconsistente}
}

\newglossaryentry{consistent}
{    name={consist\^encia},
  description={Senten\c{c}as s\~{a}o conjuntamente consistentes \ifeff{} a contradi\c{c}\~{a}o~$\ered$ \emph{n\~{a}o} pode ser provada a partir delas},
    text={consistente}
}

É fácil mostrar que algumas sentenças são inconsistentes: basta
provar a contradição~$\ered$, supondo todas as sentenças como
premissas. Mostrar que algumas sentenças \emph{não} são inconsistentes é muito
mais difícil. Isso exigiria mais do que fornecer uma ou duas provas;
exigiria mostrar que nenhuma prova de certo tipo é
\emph{possível}.

\
\\
Esta tabela resume se uma ou duas provas são suficientes ou se precisamos raciocinar sobre todas as provas possíveis.

\begin{center}
\begin{tabular}{l l l}
%\cline{2-3}
 & \textbf{Sim} & \textbf{Não}\\
 \hline
%\cline{2-3}
teorema? & uma prova & todas as provas possíveis\\
inconsistente? & uma prova & todas as provas possíveis\\
equivalente? & duas provas & todas as provas possíveis\\
consistente? & todas as provas possíveis & uma prova\\
\end{tabular}
\end{center}


\practiceproblems
\problempart
Mostre que cada uma das sentenças a seguir é um teorema:
\begin{compactlist}
\item $O \eif O$
\item $N \eor \enot N$
\item $J \eiff [J\eor (L\eand\enot L)]$
\item $((A \eif B) \eif A) \eif A$ 
\end{compactlist}

\problempart
Forneça provas para mostrar cada uma das afirmações a seguir:
\begin{compactlist}
\item $C\eif(E\eand G), \enot C \eif G \proves G$
\item $M \eand (\enot N \eif \enot M) \proves (N \eand M) \eor \enot M$
\item $(Z\eand K)\eiff(Y\eand M), D\eand(D\eif M) \proves Y\eif Z$
\item $(W \eor X) \eor (Y \eor Z), X\eif Y, \enot Z \proves W\eor Y$
\end{compactlist}

\problempart
Mostre que cada um dos pares de sentenças a seguir é interderivável:
\begin{compactlist}
\item $R \eiff E$, $E \eiff R$
\item $G$, $\enot\enot\enot\enot G$
\item $T\eif S$, $\enot S \eif \enot T$
\item $U \eif I$, $\enot(U \eand \enot I)$
\item $\enot (C \eif D), C \eand \enot D$
\item $\enot G \eiff H$, $\enot(G \eiff H)$ 
\end{compactlist}

\problempart
Se você sabe que $\metav{A}\proves\metav{B}$, o que pode dizer sobre $(\metav{A}\eand\metav{C})\proves\metav{B}$? E sobre $(\metav{A}\eor\metav{C})\proves\metav{B}$? Explique suas respostas.

\

\problempart Neste capítulo, afirmamos que é tão difícil mostrar que duas sentenças não são interderiváveis quanto mostrar que uma sentença não é um teorema. Por que fizemos essa afirmação? (\emph{Dica}: pense em uma sentença que seria um teorema \ifeff{} \metav{A} e \metav{B} fossem interderiváveis.)


\chapter{Regras derivadas}\label{s:Derived}
Nesta seção, veremos por que introduzimos as regras do nosso sistema de provas em dois grupos separados. Em particular, queremos mostrar que as regras adicionais de \cref{s:Further} não são, estritamente falando, necessárias, mas podem ser derivadas das regras básicas de \cref{s:BasicTFL}.

\section{Derivação da reiteração}
Para ilustrar o que significa derivar uma \emph{regra} a partir de outras regras, considere primeiro a reiteração. Ela é uma regra básica do nosso sistema, mas também não é necessária. Suponha que você tenha alguma sentença em alguma linha da sua dedução:
\begin{fitchproof}
	\have[m]{a}{\metav{A}}
\end{fitchproof}
Agora você quer repetir essa sentença em alguma linha $k$. Poderia simplesmente invocar a regra~R. Mas também pode fazer isso com outras regras básicas de \cref{s:BasicTFL}:
\begin{fitchproof}
	\have[m]{a}{\metav{A}}
	\have[k]{aa}{\metav{A} \eand \metav{A}}\ai{a, a}
	\have{a2}{\metav{A}}\ae{aa}
\end{fitchproof}
Para deixar claro: isto não é uma prova. É, antes, um \emph{esquema} de prova. Afinal, ele usa uma variável, $\metav{A}$, em vez de uma sentença da LVF, mas o ponto é simples: quaisquer sentenças da LVF que substituíssemos por $\metav{A}$, e quaisquer linhas em que estivéssemos trabalhando, permitiriam produzir uma prova genuína. Portanto, você pode pensar nisso como uma receita para produzir provas.

De fato, essa é uma receita que mostra o seguinte: tudo o que podemos provar usando a regra R, podemos provar (com uma linha a mais) usando apenas as regras básicas de \cref{s:BasicTFL}, sem~R. É isso que significa dizer que a regra~R pode ser derivada das outras regras básicas: tudo o que pode ser justificado usando~R pode ser justificado usando apenas as outras regras básicas.


\section{Derivação do silogismo disjuntivo}
Suponha que você esteja em uma prova e tenha algo desta forma:
\begin{fitchproof}
	\have[m]{ab}{\metav{A}\eor\metav{B}}
	\have[n]{na}{\enot \metav{A}}
\end{fitchproof}
Agora você quer provar $\metav{B}$ na linha $k$. Pode fazer isso com a regra DS, introduzida em \cref{s:Further}, mas também pode fazê-lo com as regras \emph{básicas} de \cref{s:BasicTFL}:
	\begin{fitchproof}
		\have[m]{ab}{\metav{A}\eor\metav{B}}
		\have[n]{na}{\enot \metav{A}}
		\open
			\hypo[k]{a}{\metav{A}}\AS
			\have{red}{\ered}\ne{na, a}
			\have{b1}{\metav{B}}\re{red}
		\close
		\open
			\hypo{b}{\metav{B}}\AS
			\have{b2}{\metav{B}}\by{R}{b}
		\close
	\have{con}{\metav{B}}\oe{ab, a-b1, b-b2}
\end{fitchproof}
Assim, a regra DS, novamente, pode ser derivada das nossas regras mais básicas. Adicioná-la ao nosso sistema não tornou possíveis novas provas. Sempre que usar a regra DS, você poderia acrescentar algumas linhas e provar a mesma coisa usando apenas as nossas regras básicas. Ela é uma regra \emph{derivada}.

\section{Derivação do modus tollens}
Suponha que você tenha o seguinte na sua prova:
\begin{fitchproof}
	\have[m]{ab}{\metav{A}\eif\metav{B}}
	\have[n]{a}{\enot\metav{B}}
\end{fitchproof}
Agora você quer provar $\enot \metav{A}$ na linha $k$. Pode fazer isso com a regra MT, introduzida em \cref{s:Further}. Igualmente, pode fazê-lo com as regras \emph{básicas} de \cref{s:BasicTFL}:
\begin{fitchproof}
	\have[m]{ab}{\metav{A}\eif\metav{B}}
	\have[n]{nb}{\enot\metav{B}}
		\open
		\hypo[k]{a}{\metav{A}}\AS
		\have{b}{\metav{B}}\ce{ab, a}
		\have{nb1}{\ered}\ne{nb, b}
		\close
	\have{no}{\enot\metav{A}}\ni{a-nb1}
\end{fitchproof}
Novamente, a regra MT pode ser derivada das regras \emph{básicas} de \cref{s:BasicTFL}.

\section{Derivação da eliminação da dupla negação}
Considere o seguinte esquema de dedução:
	\begin{fitchproof}
	\have[m]{m}{\enot \enot \metav{A}}
	\open
		\hypo[k]{a}{\enot\metav{A}}\AS
		\have{a1}{\ered}\ne{m, a}
	\close
	\have{con}{\metav{A}}\ip{a-a1}
\end{fitchproof}
Novamente, podemos derivar a regra DNE das regras \emph{básicas} de \cref{s:BasicTFL}.

\section{Derivação do terceiro excluído}
Suponha que você queira provar algo usando a regra LEM, isto é, que tenha na sua prova
\begin{fitchproof}
  \open
  \hypo[m]{a}{\metav{A}}\AS
  \have[n]{aaa}{\metav{B}}
  \close
  \open
  \hypo[k]{b}{\enot\metav{A}}\AS
  \have[l]{bbb}{\metav{B}}
  \close
\end{fitchproof}
Agora você quer provar $\metav{B}$ na linha $l+1$. A regra LEM de
\cref{s:Further} permitiria fazer isso. Mas você consegue fazê-lo usando
as regras \emph{básicas} de \cref{s:BasicTFL}?

Uma opção é primeiro provar $\metav{A} \eor \enot\metav{A}$ e então aplicar $\eor$E, isto é, fazer uma prova por casos:
\begin{fitchproof}
  \open
  \hypo[m]{a}{\metav{A}}\AS
  \have[n]{aaa}{\metav{B}}
  \close
  \open
  \hypo[k]{b}{\enot\metav{A}}\AS
  \have[l]{bbb}{\metav{B}}
  \close
  \ellipsesline
  \have[i]{tnd}{\metav{A} \eor \enot \metav{A}}
  \have[i+1]{fin}{\metav{B}}\oe{tnd, a-aaa,b-bbb}
\end{fitchproof}
(Demos uma prova de $\metav{A} \eor \enot\metav{A}$ usando apenas as nossas regras básicas em \cref{s:proofLEM}.)

Eis outra maneira, um pouco mais complicada que as anteriores. O que você tem de fazer é inserir as duas subprovas dentro de outra subprova. A suposição da subprova será $\enot \metav{B}$ e a última linha será $\ered$. Assim, a subprova completa é do tipo de que você precisa para concluir \metav{B} usando IP. Dentro da prova, você teria de trabalhar um pouco mais para obter~$\ered$:
\begin{fitchproof}
  \open
  \hypo[m]{nb}{\enot\metav{B}}\AS
  \open
  \hypo{a}{\metav{A}}\AS
  \ellipsesline
  \have[n]{aaa}{\metav{B}}
  \have{aaaa}{\ered}\ne{nb, aaa}
  \close
  \open
  \hypo{b}{\enot\metav{A}}\AS
  \ellipsesline
  \have[l]{bbb}{\metav{B}}
  \have{bbbb}{\ered}\ne{nb, bbb}
  \close
  \have{na}{\enot\metav{A}}\ni{(a)-(aaaa)}
  \have{nna}{\enot\enot\metav{A}}\ni{(b)-(bbbb)}
  \have{bot}{\ered}\ne{nna, na}
  \close
  \have{B}{\metav{B}}\ip{nb-(bot)}
\end{fitchproof}
Observe que, como adicionamos uma suposição no alto e conclusões adicionais dentro das subprovas, os números das linhas mudam. Talvez você precise examinar isso por algum tempo antes de entender o que está acontecendo.


\section{Derivação das regras de De Morgan}
Eis uma demonstração de como poderíamos derivar a primeira regra de De Morgan:
 	\begin{fitchproof}
		\have[m]{nab}{\enot (\metav{A} \eand \metav{B})}
		\open
			\hypo[k]{a}{\metav{A}}\AS
			\open
				\hypo{b}{\metav{B}}\AS
				\have{ab}{\metav{A} \eand \metav{B}}\ai{a,b}
				\have{nab1}{\ered}\ne{nab, ab}
			\close
			\have{nb}{\enot \metav{B}}\ni{(b)-(nab1)}
			\have{dis}{\enot\metav{A} \eor \enot \metav{B}}\oi{nb}
		\close
		\open
			\hypo{na1}{\enot \metav{A}}\AS
			\have{dis1}{\enot\metav{A} \eor \enot \metav{B}}\oi{na1}
		\close
		\have{con}{\enot \metav{A} \eor \enot \metav{B}}\tnd{a-(dis), (na1)-(dis1)}
	\end{fitchproof}
Eis uma demonstração de como poderíamos derivar a segunda regra de De Morgan:
 	\begin{fitchproof}
		\have[m]{nab}{\enot \metav{A} \eor \enot \metav{B}}
		\open
			\hypo[k]{ab}{\metav{A} \eand \metav{B}}\AS
			\have{a}{\metav{A}}\ae{ab}
			\have{b}{\metav{B}}\ae{ab}
			\open
				\hypo{na}{\enot \metav{A}}\AS
				\have{c1}{\ered}\ne{na, a}
			\close
			\open
				\hypo{nb}{\enot \metav{B}}\AS
				\have{c2}{\ered}\ne{nb, b}
			\close
			\have{con2}{\ered}\oe{nab, (na)-(c1), (nb)-(c2)}
		\close
		\have{nab}{\enot (\metav{A} \eand \metav{B})}\ni{ab-(con2)}
	\end{fitchproof}
Demonstrações semelhantes poderiam explicar como derivar a terceira e a quarta regras de De Morgan. Deixamos isso como exercício.



\practiceproblems

\problempart
Forneça esquemas de prova que justifiquem a inclusão da terceira e da quarta regras de De Morgan como regras derivadas.

\

\problempart
As provas que você ofereceu em resposta aos exercícios práticos de \crefrange{s:Further}{s:ProofTheoreticConcepts} usaram regras derivadas. Substitua o uso de regras derivadas, nessas provas, apenas por regras básicas. Você encontrará alguma ``repetição'' nas provas resultantes; nesses casos, ofereça uma prova simplificada usando apenas regras básicas. (Isso lhe dará uma noção tanto do poder das regras derivadas quanto da maneira como todas as regras interagem.)

\

\problempart
Forneça uma prova de $\metav{A} \eor \enot\metav{A}$. Em seguida,
forneça uma prova que \emph{use apenas as regras básicas}.

\

\problempart
Mostre que, se você tivesse LEM como regra básica, poderia justificar IP como regra derivada. Isto é, suponha que você tivesse a prova:
\begin{fitchproof}
  \open
  \hypo[m]{a}{\enot\metav{A}}\AS
  \have[\ ]{aa}{\dots}
  \have[n]{aaa}{\ered}
  \close
\end{fitchproof}
Como você poderia usá-la para provar $\metav{A}$ sem usar IP, mas usando LEM e todas as outras regras básicas?

\

\problempart
Forneça uma prova da primeira regra de De Morgan usando apenas as regras básicas, em particular, \emph{sem usar LEM}. (É claro que você pode combinar a prova usando LEM com a prova \emph{de}~LEM. Tente encontrar uma prova diretamente.)

\chapter{Correção e completude}
\label{sec:soundness_and_completeness}

Em \cref{s:ProofTheoreticConcepts}, vimos que podíamos usar
derivações para testar coisas semelhantes àquelas para as quais usávamos tabelas-verdade.
Podíamos não apenas mostrar que uma conclusão é derivável a partir de algumas
premissas, mas também usar derivações para testar se uma sentença é um teorema ou
se um par de sentenças é interderivável. Também começamos a usar o
turnstile simples da mesma maneira que usávamos o turnstile duplo. Se
conseguíssemos provar que $\metav{A}$ era uma tautologia usando uma tabela-verdade,
escrevíamos $\entails \metav{A}$; e, se conseguíssemos prová-la usando uma
derivação, escrevíamos $\proves \metav{A}$.

Talvez, nesse ponto, você tenha se perguntado se os dois tipos de turnstile
sempre funcionam da mesma maneira. Se você pode mostrar que $\metav{A}$ é uma
tautologia usando tabelas-verdade, também pode sempre mostrar que ela é um
teorema usando uma derivação? O inverso é verdadeiro? Essas coisas também
valem para argumentos válidos e pares de sentenças equivalentes? Acontece que
a resposta a todas essas perguntas, e a muitas outras semelhantes, é sim.
Podemos mostrar isso provando que os conceitos baseados em tabelas-verdade e
os conceitos correspondentes baseados em derivações são equivalentes. Isto é,
mostramos que, por exemplo, $\entails \metav{A}$ \ifeff{}
$\proves \metav{A}$ para qualquer sentença~$\metav{A}$.

Para começar, precisamos definir todos os nossos conceitos lógicos separadamente para tabelas-verdade e derivações. Boa parte desse trabalho já foi feita. Tratamos de todas as definições por tabela-verdade em \cref{s:SemanticConcepts}. Também já fornecemos definições prova-teóricas para teoremas e pares de sentenças interderiváveis. As outras definições seguem naturalmente. Para a maioria das propriedades lógicas, podemos conceber um teste usando derivações; aquelas que não podemos testar diretamente podem ser definidas em termos dos conceitos que conseguimos definir.

Por exemplo, definimos um teorema como uma sentença que pode ser derivada
sem quaisquer premissas (p.~\pageref{def:syntactic_tautology_in_sl}).
Podemos definir uma \define{senten\c{c}a inconsistente na LVF}%
\label{def:syntactic_contradiction_in_sl} como uma sentença a partir da qual
podemos derivar a contradição~$\ered$.\footnote{Observe que $\metav{A}
\proves \ered$ \ifeff{} $\lnot \metav{A}$ é um teorema.} A definição sintática
de uma sentença contingente é um pouco diferente. Não dispomos de nenhum
método prático e finito para provar que uma sentença é contingente usando
derivações, como fizemos usando tabelas-verdade. Portanto, temos de nos
contentar em definir negativamente ``sentença contingente''. Uma sentença é
\define{contingente prova-teoricamente na LVF}
\label{def:syntactically_contingent_in_sl} se não é nem um teorema nem uma
sentença inconsistente.

Uma coleção de sentenças é \define{inconsistente na LVF}
\label{def:syntactically_inconsistent_ in_sl} se, e somente se, podemos
derivar a contradição~$\ered$ a partir dela. A consistência, por outro lado,
é semelhante à contingência, pois não temos um método prático e finito para
testá-la diretamente. Portanto, novamente temos de definir o termo
negativamente. Uma coleção de sentenças é \define{consistente na LVF}
\label{def:syntactically consistent in SL} se, e somente se, não é
inconsistente.

Finalmente, um argumento é \define{demonstravelmente v\'alido na LVF}
\label{def:syntactically_valid_in_SL} se, e somente se, existe uma derivação
da sua conclusão a partir das suas premissas. Todas essas definições são
apresentadas em \cref{table:truth_tables_or_derivations}.

\def\tmptable{\textbf{Conceitos} 		&	\textbf{Definição por tabela-verdade (semântica)} 	&	\textbf{Definição prova-teórica (sintática)} \\ \hline \hline

Tautologia/ teorema   &	Uma sentença cuja tabela-verdade tem apenas V sob o conectivo principal & Uma sentença que pode ser derivada sem quaisquer premissas.	 \\ \hline
 
Contradição/ sentença inconsistente		&	Uma sentença cuja tabela-verdade
tem apenas F sob o conectivo principal  &	Uma sentença a partir da qual a
contradição~$\ered$ pode ser derivada\\ \hline

Sentença contingente	&	Uma sentença cuja tabela-verdade contém V e F sob o conectivo principal & Uma sentença que não é um teorema nem é inconsistente \\ \hline

Sentenças equivalentes/ interderiváveis &	As colunas sob os conectivos principais são idênticas.& As sentenças podem ser derivadas uma da outra	\\ \hline

Sentenças insatisfatíveis/ inconsistentes	&	Sentenças que não têm uma única linha na tabela-verdade em que sejam todas verdadeiras.	& Sentenças a partir das quais se pode derivar a contradição~$\ered$ \\ \hline

Sentenças satisfatíveis/ consistentes	&	Sentenças que têm pelo menos uma
linha na tabela-verdade em que são todas verdadeiras. & Sentenças a partir
das quais não se pode derivar a contradição~$\ered$	\\ \hline

Argumento válido		&	Um argumento cuja tabela-verdade não tem linhas em que haja V sob os conectivos principais das premissas e F sob o conectivo principal da conclusão.  & Um argumento cuja conclusão pode ser derivada a partir das premissas.}

\ifHTMLtarget
\begin{table}
\begin{tabularx}{\textwidth}{X||X|X}
\tmptable
\end{tabularx}
\caption{Duas maneiras de definir conceitos lógicos.}
\label{table:truth_tables_or_derivations}
\end{table}
\else
\begin{sidewaystable}\small
%\tabulinesep=1ex
%\begin{tabu}{X[.5,c,m] ||X[1,l,m] |X[1,l,m]}
\begin{tabularx}{\textwidth}{>{\hsize=.4\hsize\raggedright\arraybackslash}X||>{\raggedright\arraybackslash}X|>{\raggedright\arraybackslash}X@{}}
\tmptable
\end{tabularx}
%\end{tabu}
\caption{Duas maneiras de definir conceitos lógicos.}
\label{table:truth_tables_or_derivations}
\end{sidewaystable}
\fi

Agora temos pares correspondentes de conceitos definidos uma vez
semanticamente (usando valorações e tabelas-verdade) e outra vez
prova-teoricamente (com base na dedução natural). Como podemos
estabelecer que essas definições sempre funcionam da mesma maneira? Uma
prova completa está muito além do escopo deste livro. Contudo, podemos
esboçar como ela seria. Vamos nos concentrar em mostrar que as duas noções
de validade são equivalentes. A partir disso, os outros conceitos seguirão
rapidamente. A prova terá de seguir em duas direções. Primeiro, teremos de
mostrar que as coisas que são válidas prova-teoricamente também são
semanticamente válidas. Em outras palavras, tudo o que podemos provar
usando derivações também poderia ser provado usando tabelas-verdade. Em
termos simbólicos, queremos mostrar que $\proves$ implica
$\entails$. Depois, precisaremos mostrar a implicação na outra direção:
$\entails$ implica $\proves$.

\newglossaryentry{soundness}
{
name=correção,
description={Uma propriedade de sistemas lógicos que vale se, e somente se, $\proves $ implica $\entails $}
}

Esse argumento de $\proves $ para $\entails $ é o problema da
\define{\gls{soundness}}.\label{def:soundness} Um sistema de provas é
\define{correto} se não há derivações de argumentos que possam ser
mostrados inválidos por tabelas-verdade. \label{def_Soundness} Demonstrar
que o sistema de provas é correto exigiria mostrar que \emph{toda} prova
possível é a prova de um argumento válido. Não bastaria simplesmente ter
sucesso ao tentar provar muitos argumentos válidos e falhar ao tentar provar
argumentos inválidos. A prova em si requer algum cuidado. Em suma, ela
envolve mostrar que toda regra de inferência preserva a propriedade de que
sua conclusão é consequência lógica das premissas da prova juntamente com
quaisquer suposições abertas. Os detalhes dessa prova podem ser encontrados
em \cref{ch:Soundness}.

Correção significa que $\metav{A} \proves \metav{B}$ implica $\metav{A}
\entails \metav{B}$. E quanto à outra direção, isto é, por que pensar
que \emph{todo} argumento que pode ser mostrado válido usando tabelas-verdade
também pode ser provado usando uma derivação?

\newglossaryentry{completeness}
{
name=completude,
description={Uma propriedade de sistemas lógicos que vale se, e somente se, $\entails $ implica $\proves $}
}

Esse é o problema da completude. Um sistema de provas tem a propriedade da
\define{\gls{completeness}}\label{def:completeness} se, e somente se, existe
uma derivação de todo argumento semanticamente válido. Provar que um sistema
é completo geralmente é mais difícil do que provar que ele é correto. Provar
que um sistema é correto equivale a mostrar que todas as regras do sistema de
provas funcionam como deveriam. Mostrar que um sistema é completo significa
mostrar que incluímos \emph{todas} as regras de que precisamos, sem deixar
nenhuma de fora. Demonstrar isso está além do escopo deste livro. O ponto
importante é que, felizmente, o sistema de provas para a LVF é correto e
completo. Isso não vale para todos os sistemas de provas ou todas as
linguagens formais. Como isso vale para a LVF, podemos escolher entre
fornecer provas e fornecer tabelas-verdade---o que for mais fácil para a
tarefa em questão.

Agora que sabemos que o método das tabelas-verdade é intercambiável com o
método das derivações, podemos escolher qual método usar em cada problema.
Os estudantes frequentemente preferem usar tabelas-verdade, porque elas
podem ser produzidas de maneira puramente mecânica, o que parece ``mais
fácil''. Contudo, já vimos que as tabelas-verdade se tornam impossivelmente
grandes depois de apenas algumas letras sentenciais. Por outro lado, há
algumas situações em que simplesmente não é possível usar provas. Definimos
sintaticamente uma sentença contingente como uma sentença que não pode ser
provada como tautologia nem como contradição. Não há uma maneira prática de
provar esse tipo de afirmação negativa. Nunca saberemos se não existe alguma
prova, ainda não encontrada, de que uma sentença é uma contradição. Nessa
situação, não nos resta nada além de recorrer às tabelas-verdade. Do mesmo
modo, podemos usar derivações para provar que duas sentenças são equivalentes,
mas e se quisermos provar que elas \emph{não} são equivalentes? Não temos como
provar que nunca encontraremos a prova relevante. Portanto, precisamos
recorrer novamente às tabelas-verdade.

\Cref{table.ProofOrModel} resume quando é melhor fornecer provas e quando é
melhor fornecer tabelas-verdade.

\def\tmptable{\textbf{Propriedade lógica} 	&	\textbf{Para prová-la, apresente} 	&	\textbf{Para mostrar que ela não vale} \\ \hline \hline
Ser uma tautologia/teorema 		& Derive a sentença 						& Encontre uma linha falsa na tabela-verdade da sentença \\ \hline
Ser uma contradição 	& Derive $\ered$ a partir da sentença 		 & Encontre uma linha verdadeira na tabela-verdade da sentença\\ \hline
Contingência 			& Encontre uma linha falsa e uma linha verdadeira na tabela-verdade da sentença & Prove a sentença ou sua negação\\ \hline
Equivalência 			& Derive cada sentença a partir da outra 		 & Encontre uma linha nas tabelas-verdade das sentenças em que elas tenham valores diferentes\\ \hline
Consistência 		& Encontre uma linha na tabela-verdade da sentença em que todas sejam verdadeiras & Derive $\ered$ a partir das sentenças\\ \hline
Validade 				& Derive a conclusão a partir das premissas & Encontre uma linha na tabela-verdade em que as premissas sejam verdadeiras e a conclusão seja falsa. \\ }

\ifHTMLtarget
\begin{table}
\begin{tabularx}{\textwidth}{X||X|X}
\tmptable
\end{tabularx}
\caption{Quando fornecer uma tabela-verdade e quando fornecer uma prova.}
\label{table.ProofOrModel}
\end{table}
\else
\begin{table}\small
%\tabulinesep=1ex
%\begin{tabu}{X[.7,c,m] ||X[1,l,m] |X[1,l,m]}
\begin{tabularx}{\textwidth}{>{\hsize=.7\hsize\raggedright\arraybackslash}X||X|X}
\tmptable
\end{tabularx}
%\end{tabu}
\caption{Quando fornecer uma tabela-verdade e quando fornecer uma prova.}
\label{table.ProofOrModel}
\end{table}
\fi



\practiceproblems
\noindent\problempart Use uma derivação ou uma tabela-verdade para cada item a seguir.
\begin{compactlist}%[label=(\arabic*)]
\item Mostre que $A \eif [((B \eand C) \eor D) \eif A]$ é um teorema.
\item Mostre que $A \eif (A \eif B)$ não é um teorema.
\item Mostre que a sentença $A \eif \enot{A}$ não é uma contradição.
\item Mostre que a sentença $A \eiff \enot A$ é uma contradição.
\item Mostre que a sentença $ \enot (W \eif (J \eor J)) $ é contingente.
\item Mostre que a sentença $ \enot(X \eor (Y \eor Z)) \eor (X \eor (Y \eor Z))$ não é contingente.
 \item Mostre que a sentença $B \eif \enot S$ é equivalente à sentença $\enot \enot B \eif \enot S$.
\item Mostre que a sentença $ \enot (X \eor O) $ não é equivalente à sentença $X \eand O$.
\item Mostre que as sentenças $\enot(A \eor B)$, $C$, $C \eif A$ são conjuntamente inconsistentes.
\item Mostre que as sentenças $\enot(A \eor B)$, $\enot{B}$, $B \eif A$ são conjuntamente consistentes.
\item Mostre que $\enot(A \eor (B \eor C)) $ \therefore $ \enot{C}$ é válido.
\item Mostre que $\enot(A \eand (B \eor C))$ \therefore $ \enot{C}$ é inválido.
\end{compactlist}


\noindent\problempart Use uma derivação ou uma tabela-verdade para cada item a seguir.
\begin{compactlist}%[label=(\arabic*)]
\item Mostre que $A \eif (B \eif A)$ é um teorema.
\item Mostre que $\enot (((N \eiff Q) \eor Q) \eor N)$ não é um teorema.
\item Mostre que $ Z \eor (\enot Z \eiff Z) $ é contingente.
\item Mostre que $ (L \eiff ((N \eif N) \eif L)) \eor H $ não é contingente.
\item Mostre que $ (A \eiff A) \eand (B \eand \enot B)$ é uma contradição.
\item Mostre que $ (B \eiff (C \eor B)) $ não é uma contradição.
\item Mostre que $ ((\enot X \eiff X) \eor X) $ é equivalente a $X$.
\item Mostre que $F \eand (K \eand R) $ não é equivalente a $ (F \eiff (K \eiff R)) $.
\item Mostre que as sentenças $ \enot (W \eif W)$, $(W \eiff W) \eand W$, $E \eor (W \eif \enot (E \eand W))$ são conjuntamente inconsistentes.
\item Mostre que as sentenças $\enot R \eor C $, $(C \eand R) \eif \enot R$, $(\enot (R \eor C) \eif R) $ são conjuntamente consistentes.
\item Mostre que $\enot \enot (C \eiff \enot C), ((G \eor C) \eor G) \therefore ((G \eif C) \eand G) $ é válido.
\item Mostre que $ \enot \enot L,  (C \eif \enot L) \eif C) \therefore \enot C$ é inválido.
\end{compactlist}
